\documentclass[12pt,a4paper,oneside]{book}

\usepackage[utf8]{inputenc}
\usepackage[T1]{fontenc}
\usepackage[brazil]{babel}
\usepackage{lmodern}
\usepackage{amsmath,amssymb,amsfonts,amsthm}
\usepackage{mathtools}
\usepackage{bm}
\usepackage{geometry}
\usepackage{xcolor}
\usepackage{booktabs}
\usepackage{array}
\usepackage{longtable}
\usepackage{enumitem}
\usepackage{hyperref}
\usepackage{cleveref}
\usepackage[most]{tcolorbox}
\usepackage{fancyhdr}
\usepackage{microtype}
\usepackage{setspace}

\geometry{top=3.2cm,bottom=3.2cm,left=3.1cm,right=3.1cm,headheight=15pt}
\doublespacing
\setlength{\parskip}{3pt}
\setlist[itemize]{topsep=3pt,itemsep=2pt}
\setlist[enumerate]{topsep=3pt,itemsep=2pt}

\hypersetup{
  colorlinks=true,
  linkcolor=blue!55!black,
  citecolor=green!45!black,
  urlcolor=cyan!55!black,
  pdftitle={Relatividade Geral: Apostila para 72 Aulas},
  pdfauthor={Rafael -- UDESC/CCT}
}

\pagestyle{fancy}
\fancyhf{}
\fancyhead[L]{\small Relatividade Geral}
\fancyhead[R]{\small \leftmark}
\fancyfoot[C]{\thepage}

\theoremstyle{definition}
\newtheorem{definition}{Definição}[chapter]
\newtheorem{example}{Exemplo}[chapter]
\newtheorem{exercise}{Exercício}[chapter]
\newtheorem{activity}{Atividade}[chapter]

\theoremstyle{plain}
\newtheorem{proposition}{Proposição}[chapter]

\theoremstyle{remark}
\newtheorem{remark}{Observação}[chapter]

\newtcolorbox{objetivos}{
  colback=blue!4!white,
  colframe=blue!55!black,
  title=Objetivos,
  fonttitle=\bfseries,
  breakable
}

\newtcolorbox{resultado}{
  colback=green!4!white,
  colframe=green!45!black,
  title=Resultado central,
  fonttitle=\bfseries,
  breakable
}

\newtcolorbox{roteiro}{
  colback=gray!5!white,
  colframe=gray!65!black,
  title=Roteiro de aula,
  fonttitle=\bfseries,
  breakable
}

\newtcolorbox{leitura}{
  colback=yellow!7!white,
  colframe=orange!70!black,
  title=Leitura orientada,
  fonttitle=\bfseries,
  breakable
}

\newcommand{\dd}{\mathrm{d}}
\newcommand{\D}{\mathrm{D}}
\newcommand{\pd}{\partial}
\newcommand{\Lie}{\mathcal{L}}
\newcommand{\Riem}{R^{\rho}{}_{\sigma\mu\nu}}
\newcommand{\Ric}{R_{\mu\nu}}
\newcommand{\Ein}{G_{\mu\nu}}
\newcommand{\T}{T_{\mu\nu}}
\newcommand{\M}{\mathcal{M}}
\newcommand{\Lag}{\mathcal{L}}
\newcommand{\eps}{\varepsilon}
\newcommand{\half}{\frac{1}{2}}
\newcommand{\etaab}{\eta_{\alpha\beta}}
\newcommand{\munu}{\mu\nu}
\begin{document}

\frontmatter

\begin{titlepage}
\centering
\vspace*{2cm}
{\Large Universidade do Estado de Santa Catarina\\
Centro de Ciências Tecnológicas}\\[2.5cm]

{\Huge\bfseries Relatividade Geral}\\[0.5cm]
{\LARGE Apostila para um curso de 72 encontros}\\[0.8cm]
{\Large Aulas de 50 minutos}\\[2cm]

{\large Material original de apoio didático}\\[0.3cm]
{\large Referência principal: Bernard Schutz, \textit{A First Course in General Relativity}, 3ª ed.}\\[2.5cm]

{\large 2026}
\vfill
{\small Esta apostila foi preparada como material de ensino. Ela usa o livro de Bernard Schutz como guia de organização e referência bibliográfica, mas apresenta texto, exemplos e exercícios em formulação própria.}
\end{titlepage}

\chapter*{Prefácio}
\addcontentsline{toc}{chapter}{Prefácio}

Estas notas foram organizadas para uma disciplina de Relatividade Geral com
72 encontros de 50 minutos. O objetivo é conduzir o estudante desde a
relatividade especial e a linguagem tensorial até aplicações centrais da teoria:
ondas gravitacionais, estrelas relativísticas, buracos negros e cosmologia.

O livro de Bernard Schutz, \textit{A First Course in General Relativity}, terceira
edição, foi usado como referência principal para a sequência conceitual. A
apostila não substitui o livro: ela funciona como roteiro de aula, resumo
matemático, coleção de atividades e ponte entre o quadro-negro e os problemas
que o estudante deve resolver por conta própria.

O curso assume familiaridade com cálculo vetorial, álgebra linear, equações
diferenciais ordinárias, mecânica clássica e eletromagnetismo em nível de
graduação. Quando uma ferramenta matemática aparece pela primeira vez, ela é
introduzida operacionalmente: primeiro se explica para que serve, depois se
formaliza.

\tableofcontents

\mainmatter

\chapter{Plano do curso}

\section{Ementa expandida}

Relatividade restrita. Espaço-tempo de Minkowski. Diagramas espaço-tempo.
Quadrivetores. Energia, momento e fótons. Tensores, formas lineares, métrica e
notação de índices. Fluido perfeito e tensor energia-momento. Princípio da
equivalência. Coordenadas curvilíneas. Derivada covariante, transporte paralelo,
geodésicas e curvatura. Equações de Einstein. Limite newtoniano. Ondas
gravitacionais. Soluções esfericamente simétricas. Estrelas relativísticas.
Geometria de Schwarzschild. Buracos negros. Cosmologia de Friedmann-Lemaître-
Robertson-Walker. Problemas selecionados de astrofísica relativística.

\section{Objetivos de aprendizagem}

Ao final do curso, o estudante deverá ser capaz de:

\begin{itemize}
  \item usar quadrivetores e tensores em espaço-tempo plano e curvo;
  \item interpretar a métrica como objeto geométrico e físico;
  \item calcular símbolos de Christoffel, geodésicas e curvatura em exemplos simples;
  \item explicar o princípio da equivalência e sua relação com a geometrização da gravidade;
  \item derivar o limite newtoniano das equações de Einstein;
  \item aplicar a métrica de Schwarzschild a órbitas, desvio da luz, redshift gravitacional e horizonte;
  \item descrever a física básica de ondas gravitacionais e seus detectores;
  \item formular as equações cosmológicas de Friedmann e interpretar soluções simples.
\end{itemize}

\section{Distribuição dos 72 encontros}

\begin{center}
\small
\begin{tabular}{cp{3.1cm}p{8.1cm}}
\toprule
Encontros & Módulo & Conteúdo \\
\midrule
1--6 & Relatividade especial & Postulados, intervalos, diagramas, transformações de Lorentz, paradoxo dos gêmeos. \\
7--12 & Vetores relativísticos & Quadrivelocidade, quadrimomento, produto escalar, fótons, colisões e decaimentos. \\
13--18 & Tensores em Minkowski & Um-formas, métrica, tensores mistos, mudança de base, derivadas. \\
19--22 & Fluidos relativísticos & Fluxo de partículas, tensor energia-momento, fluido perfeito, conservação. \\
23--30 & Curvatura em coordenadas & Coordenadas polares, bases não coordenadas, símbolos de Christoffel, geodésicas. \\
31--38 & Variedades curvas & Derivada covariante, transporte paralelo, tensor de Riemann, Ricci, escalar de curvatura, Bianchi. \\
39--44 & Física em espaço-tempo curvo & Princípio da equivalência, referenciais locais, constantes de movimento, partículas e luz. \\
45--50 & Equações de Einstein & Motivação, ação gravitacional, limite fraco, limite newtoniano, constante cosmológica. \\
51--56 & Ondas gravitacionais & Perturbações lineares, gauge TT, polarizações, resposta de detectores, emissão quadrupolar. \\
57--62 & Estrelas e Schwarzschild & Simetria esférica, equações TOV, exterior de Schwarzschild, órbitas, lentes, precessão. \\
63--67 & Buracos negros & Horizonte, coordenadas regulares, geodésicas radiais, Kerr qualitativo, termodinâmica. \\
68--72 & Cosmologia & Métrica FLRW, redshift cosmológico, Friedmann, conteúdo material, história térmica e problemas abertos. \\
\bottomrule
\end{tabular}
\end{center}

\section{Avaliação sugerida}

\begin{itemize}
  \item Listas quinzenais de problemas: 40\%.
  \item Duas provas escritas: 40\%.
  \item Seminário final ou projeto computacional curto: 20\%.
\end{itemize}

\section{Calendário aula a aula}

\scriptsize
\begin{longtable}{cp{4.2cm}p{8.0cm}}
\toprule
Aula & Foco & Seções da apostila a cobrir \\
\midrule
\endfirsthead
\toprule
Aula & Foco & Seções da apostila a cobrir \\
\midrule
\endhead
1 & Apresentação e escala física da RG & Plano do curso: Ementa expandida, Objetivos de aprendizagem, Distribuição dos 72 encontros, Avaliação sugerida. Leitura orientadora: Prefácio. \\
2 & Eventos e intervalo & Relatividade especial: Eventos, observadores e intervalo. Iniciar Leitura física do intervalo. \\
3 & Tempo próprio e causalidade & Relatividade especial: Tempo próprio; concluir Leitura física do intervalo. \\
4 & Lorentz I & Relatividade especial: Transformações de Lorentz; exemplos de simultaneidade em Exemplos resolvidos adicionais. \\
5 & Lorentz II e contração & Relatividade especial: Contração e dilatação sem mistério; Problemas orientados, itens 1--3. \\
6 & Paradoxo dos gêmeos e fechamento de RE & Relatividade especial: Paradoxo dos gêmeos; Problemas orientados, itens 4--6; Exercícios. \\
7 & Quadrivelocidade & Vetores e momento relativístico: Quadrivelocidade; conectar com Relatividade especial: Tempo próprio. \\
8 & Quadrimomento & Vetores e momento relativístico: Quadrimomento; Energia como componente temporal. \\
9 & Conservação relativística & Vetores e momento relativístico: Conservação; exemplos de decaimento em dois corpos. \\
10 & Fótons e Doppler & Vetores e momento relativístico: Partículas massivas e partículas sem massa; Problemas orientados, itens 4--5. \\
11 & Centro de momento & Vetores e momento relativístico: Centro de momento; Problemas orientados, itens 2--3 e 6. \\
12 & Consolidação de RE e quadrivetores & Vetores e momento relativístico: Exercícios. Revisão conjunta dos capítulos Relatividade especial e Vetores e momento relativístico. \\
13 & Um-formas e contrações & Tensores no espaço-tempo plano: Vetores, um-formas e contrações. Apoio: Apêndice~\ref{ap:indices}, A filosofia dos índices e Delta de Kronecker. \\
14 & Métrica e índices & Tensores no espaço-tempo plano: Métrica; A métrica como instrumento de medida. Apoio: Apêndice~\ref{ap:indices}, Tradução para relatividade. \\
15 & Regras de índices & Tensores no espaço-tempo plano: Regra prática de índices; O que significa uma equação tensorial. Apoio: Apêndice~\ref{ap:indices}, Índices mudos e troca de nomes. \\
16 & Mudança de base & Tensores no espaço-tempo plano: Mudança de base e transformação de componentes; Exemplos resolvidos adicionais. \\
17 & Erros comuns e prática & Tensores no espaço-tempo plano: Erros comuns com índices; Problemas orientados. Apoio: Apêndice~\ref{ap:indices}, Checklist para não errar. \\
18 & Oficina de tensores planos & Tensores no espaço-tempo plano: Exercícios. Apoio opcional: Apêndice~\ref{ap:indices}, Símbolo de Levi-Civita e A identidade fundamental. \\
19 & Poeira e fluxo de partículas & Fluidos relativísticos: Poeira; Como ler o tensor energia-momento. \\
20 & Fluido perfeito & Fluidos relativísticos: Fluido perfeito; Equações de estado. \\
21 & Conservação local & Fluidos relativísticos: Conservação local e leis integrais; Problemas orientados, itens 1--4. \\
22 & Fechamento de fluidos & Fluidos relativísticos: Problemas orientados, itens 5--6; Exercícios. \\
23 & Coordenadas curvilíneas & Da gravidade à curvatura: Coordenadas curvilíneas; Christoffel não é força. \\
24 & Geodésicas I & Da gravidade à curvatura: Equação geodésica; Derivação por princípio variacional. \\
25 & Equivalência e marés & Da gravidade à curvatura: O princípio da equivalência; Marés como sinal de curvatura. \\
26 & Geodésicas II & Da gravidade à curvatura: Problemas orientados, itens 1--4. \\
27 & Esfera e exemplos 2D & Da gravidade à curvatura: Problemas orientados, itens 5--6; Exercícios. \\
28 & Variedades & Variedades, derivada covariante e curvatura: Variedades e métricas. \\
29 & Derivada covariante & Variedades, derivada covariante e curvatura: Derivada covariante; Compatibilidade métrica. \\
30 & Transporte paralelo e curvatura & Variedades, derivada covariante e curvatura: Curvatura; Curvatura intrínseca e extrínseca. \\
31 & Riemann e simetrias & Variedades, derivada covariante e curvatura: Simetrias do tensor de Riemann; Exemplos resolvidos adicionais. \\
32 & Desvio geodésico & Variedades, derivada covariante e curvatura: Desvio geodésico; Problemas orientados, itens 1--3. \\
33 & Ricci, escalar e Einstein & Variedades, derivada covariante e curvatura: Problemas orientados, itens 4--6; Resumo de fórmulas: Geometria. \\
34 & Oficina de geometria diferencial & Variedades, derivada covariante e curvatura: Exercícios. \\
35 & Movimento livre & Física em espaço-tempo curvo: Movimento livre; Energia medida por um observador. \\
36 & Observadores locais & Física em espaço-tempo curvo: Observadores não são coordenadas; Tetradas locais. \\
37 & Simetrias e Killing & Física em espaço-tempo curvo: Simetrias e quantidades conservadas; Constantes de movimento e energia local. \\
38 & Oficina de física em espaço-tempo curvo & Física em espaço-tempo curvo: Problemas orientados; Exercícios. \\
39 & Equações de campo I & Equações de Einstein: A equação de campo; Por que o tensor de Einstein. \\
40 & Equações de campo II & Equações de Einstein: Traço das equações de campo; Conservação local. \\
41 & Limite newtoniano I & Equações de Einstein: Limite newtoniano; O limite fraco com mais detalhe. \\
42 & Constante cosmológica & Equações de Einstein: Constante cosmológica como fluido; Problemas orientados, itens 1--3. \\
43 & Vácuo e ordens de grandeza & Equações de Einstein: Exemplos resolvidos adicionais; Problemas orientados, itens 4--6. \\
44 & Fechamento de Einstein & Equações de Einstein: Exercícios; Resumo de fórmulas: Einstein. \\
45 & Ondas gravitacionais I & Ondas gravitacionais: Campo fraco; Gauge e graus de liberdade. \\
46 & Polarizações & Ondas gravitacionais: Polarizações; Efeito sobre partículas teste. \\
47 & Detectores & Ondas gravitacionais: Detecção; Escalas de detecção. \\
48 & Emissão & Ondas gravitacionais: Emissão; Binárias compactas. \\
49 & Oficina de ondas & Ondas gravitacionais: Problemas orientados; Erros comuns e interpretação. \\
50 & Fechamento de ondas & Ondas gravitacionais: Exercícios. Preparar transição para soluções exatas. \\
51 & Schwarzschild I & Soluções esféricas, estrelas e Schwarzschild: Métrica de Schwarzschild; Coordenada areal. \\
52 & Redshift & Soluções esféricas, estrelas e Schwarzschild: Redshift gravitacional; Testes clássicos, parte de redshift. \\
53 & Geodésicas em Schwarzschild & Soluções esféricas, estrelas e Schwarzschild: Geodésicas e potencial efetivo. \\
54 & Órbitas circulares e fótons & Soluções esféricas, estrelas e Schwarzschild: Órbitas circulares; Órbita de fótons. \\
55 & Estrelas relativísticas & Soluções esféricas, estrelas e Schwarzschild: Estrelas relativísticas. \\
56 & Testes clássicos e exercícios & Soluções esféricas, estrelas e Schwarzschild: Testes clássicos; Problemas orientados; Exercícios. \\
57 & Horizonte & Buracos negros: Horizonte; Horizon versus singularidade. \\
58 & Coordenadas regulares & Buracos negros: Coordenadas regulares; Diagramas causais. \\
59 & Buracos negros reais & Buracos negros: Buracos negros em rotação; Buracos negros reais. \\
60 & Termodinâmica & Buracos negros: Termodinâmica; Leis da mecânica de buracos negros. \\
61 & Oficina de buracos negros & Buracos negros: Problemas orientados; Erros comuns e interpretação. \\
62 & Fechamento de buracos negros & Buracos negros: Exercícios; Resumo de fórmulas: Schwarzschild. \\
63 & Cosmologia I & Cosmologia relativística: Princípio cosmológico; Redshift cosmológico. \\
64 & Friedmann I & Cosmologia relativística: Equações de Friedmann; Eras de dominação. \\
65 & Distâncias e horizontes & Cosmologia relativística: Distâncias em cosmologia; Horizontes cosmológicos. \\
66 & Parâmetros e problemas abertos & Cosmologia relativística: Parâmetros observacionais; Problemas conceituais. \\
67 & Oficina de cosmologia & Cosmologia relativística: Problemas orientados; Exercícios, itens de matéria e radiação. \\
68 & Síntese cosmológica & Cosmologia relativística: Erros comuns e interpretação; Síntese conceitual do curso; Resumo de fórmulas: Cosmologia. \\
69 & Revisão técnica & Apêndice~\ref{ap:indices}: Checklist para não errar; Revisão de tensores, geodésicas, curvatura e Einstein. \\
70 & Projetos I & Projetos finais: Projeto 1 e Projeto 2. Definir grupos, escopo, equações e entregáveis. \\
71 & Projetos II & Projetos finais: Projeto 3 e Projeto 4. Discussão de formas de onda, Friedmann numérico e critérios de avaliação. \\
72 & Síntese final & Revisão global: Síntese conceitual do curso; Bibliografia comentada; fechamento dos projetos e preparação para avaliação final. \\
\bottomrule
\end{longtable}
\normalsize

\chapter{Relatividade especial}

\begin{objetivos}
Entender que a relatividade especial é a geometria do espaço-tempo plano.
O estudante deve sair deste módulo vendo o intervalo como a grandeza
invariante fundamental e o tempo próprio como quantidade física medida por
um relógio ao longo de uma linha de mundo.
\end{objetivos}

\section{Eventos, observadores e intervalo}

Um evento é algo que ocorre em um ponto do espaço e em um instante. Em um
referencial inercial, escrevemos suas coordenadas como
\[
  x^\mu=(t,x,y,z),
\]
onde adotaremos frequentemente unidades com $c=1$. O intervalo entre dois
eventos próximos é
\[
  \dd s^2=-\dd t^2+\dd x^2+\dd y^2+\dd z^2.
\]
A convenção de sinais usada aqui é $(-,+,+,+)$. O mesmo conteúdo físico pode
ser escrito com a convenção oposta, desde que ela seja usada de modo
consistente.

\begin{resultado}
O intervalo $\dd s^2$ é invariante sob transformações de Lorentz. Essa
invariância substitui a noção newtoniana de tempo absoluto.
\end{resultado}

Quando $\dd s^2<0$, a separação é temporal; quando $\dd s^2=0$, é nula; quando
$\dd s^2>0$, é espacial. Partículas massivas percorrem curvas temporais. A luz
percorre curvas nulas.

\section{Tempo próprio}

Para uma partícula massiva, define-se o tempo próprio $\tau$ por
\[
  \dd \tau^2=-\dd s^2.
\]
Se a partícula tem velocidade ordinária $v$, então
\[
  \dd \tau=\dd t\sqrt{1-v^2}=\frac{\dd t}{\gamma},
  \qquad
  \gamma=\frac{1}{\sqrt{1-v^2}}.
\]
Assim, o relógio comóvel à partícula mede $\tau$, não $t$.

\section{Transformações de Lorentz}

Para dois referenciais com velocidade relativa $v$ ao longo de $x$, a
transformação de Lorentz é
\[
  t'=\gamma(t-vx), \qquad
  x'=\gamma(x-vt), \qquad
  y'=y,\qquad z'=z.
\]
Ela preserva o intervalo e mistura tempo com espaço. A relatividade da
simultaneidade nasce desse acoplamento.

\begin{activity}
Mostre que dois eventos simultâneos em $S$, separados por $\Delta x\ne0$, não
são simultâneos em $S'$ se $v\ne0$. Interprete o resultado em um diagrama
espaço-tempo.
\end{activity}

\section{Paradoxo dos gêmeos}

O paradoxo dos gêmeos não é uma contradição porque as duas linhas de mundo
entre os eventos ``partida'' e ``reencontro'' não são equivalentes. O tempo
próprio acumulado é
\[
  \tau=\int \sqrt{1-v(t)^2}\,\dd t.
\]
Entre dois eventos fixos em espaço-tempo plano, a linha de mundo inercial
maximiza o tempo próprio. O gêmeo que acelera para retornar percorre uma linha
de mundo com menor tempo próprio.

\section{Leitura física do intervalo}

O intervalo é a primeira grande mudança conceitual do curso. Em mecânica
newtoniana, dois observadores podem discordar das coordenadas espaciais de um
evento, mas concordam sobre o intervalo de tempo entre eventos. Na
relatividade especial, essa separação deixa de ser absoluta. O que permanece
absoluto é uma combinação de tempo e espaço. A geometria que antes era
euclidiana no espaço tridimensional passa a ser pseudo-euclidiana no
espaço-tempo.

É útil insistir em uma frase: o sinal de $\Delta s^2$ tem conteúdo causal.
Se $\Delta s^2<0$, existe um referencial no qual os dois eventos ocorrem no
mesmo ponto espacial. Esse é o referencial momentaneamente comóvel a uma
partícula massiva que poderia ligar os eventos. Se $\Delta s^2>0$, existe um
referencial no qual os eventos são simultâneos. Nesse caso nenhum sinal causal,
limitado pela velocidade da luz, pode conectar os eventos. Se $\Delta s^2=0$,
os eventos podem ser ligados por um raio luminoso.

\begin{example}
Considere dois eventos separados por $\Delta t=5$ e $\Delta x=3$, em unidades
com $c=1$. O intervalo é
\[
  \Delta s^2=-25+9=-16.
\]
A separação é temporal. O tempo próprio entre os eventos, ao longo da reta
inercial que os conecta, é $\Delta\tau=4$. Esse número é o tempo medido por um
relógio que passa pelos dois eventos sem aceleração.
\end{example}

\begin{example}
Agora tome $\Delta t=2$ e $\Delta x=5$. Temos
\[
  \Delta s^2=-4+25=21.
\]
A separação é espacial. Não faz sentido perguntar qual partícula massiva liga
esses eventos, pois isso exigiria velocidade maior que a da luz. Contudo, faz
sentido procurar um referencial no qual os eventos sejam simultâneos.
\end{example}

\section{Contração e dilatação sem mistério}

A dilatação temporal não significa que ``o tempo realmente fica lento'' em
sentido absoluto. Significa que, entre dois eventos sobre a linha de mundo de
um relógio, o tempo próprio daquele relógio é menor que a diferença de tempo
coordenado atribuída por um observador que vê o relógio se mover. A contração
de comprimentos também depende de simultaneidade: medir o comprimento de uma
barra em movimento exige registrar as posições de suas extremidades no mesmo
instante do referencial do observador.

Essa observação elimina muitos falsos paradoxos. Quando dois observadores em
movimento relativo se descrevem mutuamente como tendo relógios dilatados, eles
não estão fazendo afirmações contraditórias sobre um mesmo par de eventos. Cada
um compara eventos diferentes, definidos por sua própria noção de
simultaneidade. A geometria de Minkowski organiza essas comparações sem
privilegiar um observador.

\begin{activity}
Desenhe uma barra em repouso em $S'$ e represente as linhas de mundo de suas
extremidades. Marque dois eventos simultâneos em $S$ e use o diagrama para
explicar por que o comprimento medido em $S$ é menor que o comprimento próprio.
\end{activity}

\section{Problemas orientados}

\begin{enumerate}
  \item \textbf{Relógio de luz.} Monte o argumento da dilatação temporal usando
  dois espelhos separados por uma distância $L$. Primeiro calcule o período no
  referencial do relógio. Depois calcule o caminho diagonal percorrido pela luz
  quando o relógio se move com velocidade $v$. A pista é aplicar Pitágoras ao
  triângulo formado por metade do percurso.

  \item \textbf{Comprimento próprio.} Uma barra está em repouso em $S'$ com
  comprimento $L_0$. Use transformações de Lorentz para mostrar que eventos
  simultâneos em $S$ nas extremidades da barra não são simultâneos em $S'$.
  Conclua que $L=L_0/\gamma$.

  \item \textbf{Ordem temporal.} Escolha dois eventos com separação espacial.
  Encontre uma velocidade $v$ tal que a ordem temporal dos eventos seja
  invertida. Discuta por que isso não viola causalidade.

  \item \textbf{Tempo próprio máximo.} Compare uma linha de mundo inercial com
  uma trajetória formada por dois trechos retilíneos de velocidades opostas.
  Calcule o tempo próprio total e interprete a desigualdade resultante.

  \item \textbf{Adição de velocidades.} Derive a fórmula
  \[
    u'=\frac{u-v}{1-uv}
  \]
  a partir das transformações de Lorentz. Verifique explicitamente que, se
  $u=1$, então $u'=1$.

  \item \textbf{Rapidez.} Defina a rapidez $\varphi$ por $v=\tanh\varphi$.
  Mostre que transformações de Lorentz sucessivas na mesma direção somam
  rapidezes. Essa variável revela a estrutura hiperbólica da relatividade.
\end{enumerate}

\section{Exemplos resolvidos adicionais}

\begin{example}
\textbf{Simultaneidade relativa.} Dois eventos ocorrem em $S$ com
$\Delta t=0$ e $\Delta x=L$. Em $S'$, que se move com velocidade $v$ em relação
a $S$,
\[
  \Delta t'=\gamma(\Delta t-v\Delta x)=-\gamma vL.
\]
Logo os eventos não são simultâneos em $S'$. O sinal negativo indica que, para
$v>0$, o evento com maior $x$ ocorre antes em $S'$. Essa conta simples é a
origem operacional da contração de comprimentos: medir uma barra em movimento
requer escolher eventos simultâneos no referencial de medida.
\end{example}

\begin{example}
\textbf{Invariância do cone de luz.} Se $\Delta x=\Delta t$ em uma dimensão
espacial, então a separação é nula. Aplicando Lorentz,
\[
  \Delta t'=\gamma(\Delta t-v\Delta x)=\gamma\Delta t(1-v),
\]
\[
  \Delta x'=\gamma(\Delta x-v\Delta t)=\gamma\Delta t(1-v).
\]
Portanto $\Delta x'=\Delta t'$. Um raio de luz continua sendo raio de luz para
qualquer observador inercial.
\end{example}

\begin{example}
\textbf{Rapidez e composição.} Escreva $v=\tanh\varphi$. Então
$\gamma=\cosh\varphi$ e $\gamma v=\sinh\varphi$. A transformação de Lorentz em
$1+1$ dimensões fica
\[
  t' = t\cosh\varphi-x\sinh\varphi,\qquad
  x' = x\cosh\varphi-t\sinh\varphi.
\]
Ela se parece com uma rotação hiperbólica. Duas transformações sucessivas
somam os parâmetros $\varphi$, o que explica por que a lei de adição de
velocidades não é simplesmente soma de velocidades.
\end{example}

\begin{example}
\textbf{Tempo próprio em trajetória por partes.} Uma nave afasta-se da Terra
com velocidade $v$ durante tempo coordenado $T/2$ e retorna com velocidade
$-v$ durante $T/2$. O tempo próprio acumulado é
\[
  \tau=T\sqrt{1-v^2}.
\]
O relógio terrestre, aproximadamente inercial nesse modelo idealizado, acumula
$T$. A diferença não é causada apenas pelo instante de aceleração, mas pela
geometria total da linha de mundo escolhida.
\end{example}

\clearpage
\section{Exercícios}

\begin{exercise}
Verifique explicitamente que a transformação de Lorentz preserva
$-\Delta t^2+\Delta x^2$.
\end{exercise}

\begin{exercise}
Uma partícula viaja com velocidade constante $v=0{,}8$ durante 10 anos no
referencial da Terra. Quanto tempo próprio passa para a partícula?
\end{exercise}

\begin{exercise}
Desenhe os cones de luz de um evento e marque regiões causalmente conectadas
ao evento.
\end{exercise}

\chapter{Vetores e momento relativístico}

\begin{objetivos}
Construir a linguagem de quadrivetores e usá-la para resolver problemas de
cinemática e dinâmica relativística sem depender de componentes particulares.
\end{objetivos}

\section{Quadrivelocidade}

A linha de mundo de uma partícula é parametrizada por $x^\mu(\tau)$. A
quadrivelocidade é
\[
  U^\mu=\frac{\dd x^\mu}{\dd\tau}.
\]
Em componentes inerciais,
\[
  U^\mu=\gamma(1,\bm v).
\]
Seu produto escalar consigo mesma é fixo:
\[
  U^\mu U_\mu=-1.
\]
Isso mostra que a quadrivelocidade não é uma velocidade comum em quatro
dimensões; ela é o vetor tangente normalizado à linha de mundo.

\section{Quadrimomento}

Para massa de repouso $m$,
\[
  p^\mu=mU^\mu=(E,\bm p),
  \qquad
  E=\gamma m,\quad \bm p=\gamma m\bm v.
\]
A relação de dispersão é
\[
  p^\mu p_\mu=-m^2
  \quad\Longrightarrow\quad
  E^2=\bm p^2+m^2.
\]
Para fótons, $m=0$ e $E=|\bm p|$.

\section{Conservação}

Em uma interação local no espaço-tempo plano,
\[
  \sum_{\rm inicial}p^\mu=\sum_{\rm final}p^\mu.
\]
Essa equação compacta contém conservação de energia e de momento. A vantagem
da forma tensorial é que, se ela é verdadeira em um referencial inercial, é
verdadeira em todos.

\begin{example}
Considere o decaimento de uma partícula de massa $M$ em duas partículas
idênticas de massa $m$, no referencial de repouso da partícula inicial. Por
simetria, os momentos finais têm módulos iguais e direções opostas. A energia
de cada produto é $E=M/2$, logo
\[
  |\bm p|=\sqrt{\frac{M^2}{4}-m^2}.
\]
O decaimento só é possível se $M\ge 2m$.
\end{example}

\section{Energia como componente temporal}

No formalismo relativístico, energia não é uma quantidade isolada: ela é a
componente temporal do quadrimomento. Isso ajuda a entender por que energia e
momento se misturam sob transformações de Lorentz. Assim como tempo e espaço
são partes de um mesmo objeto geométrico, energia e momento também são partes
de um mesmo objeto. Uma colisão que parece envolver apenas energia em um
referencial pode envolver energia e momento de modo inseparável em outro.

A massa de repouso aparece como norma do quadrimomento:
\[
  m^2=-p^\mu p_\mu.
\]
Essa equação é mais fundamental que $E=\gamma m$, pois não depende do
referencial escolhido. Em particular, a energia de uma partícula depende do
observador, mas a massa de repouso não. Para um sistema de várias partículas, a
massa invariante total é obtida a partir do quadrimomento total,
\[
  M_{\rm sist}^2=-P^\mu P_\mu,
  \qquad
  P^\mu=\sum_a p_a^\mu.
\]
Essa massa inclui energia cinética interna. Por isso um gás quente confinado em
uma caixa tem massa total ligeiramente maior que o mesmo gás frio, se o número
de partículas for o mesmo.

\begin{example}
Duas partículas sem massa, cada uma com energia $E$, movem-se em direções
opostas. Seus quadrimomentos podem ser escritos como
\[
  p_1^\mu=(E,E,0,0),\qquad p_2^\mu=(E,-E,0,0).
\]
O quadrimomento total é $P^\mu=(2E,0,0,0)$, logo
\[
  M_{\rm sist}=2E.
\]
Embora cada partícula tenha massa zero, o sistema tem massa invariante não
nula.
\end{example}

\section{Centro de momento}

O referencial do centro de momento é definido por $\bm P=0$. Ele simplifica
problemas de colisão porque toda a energia total disponível aparece como massa
invariante do sistema. Em uma colisão de duas partículas, a quantidade
\[
  s=-(p_1+p_2)^2
\]
é invariante. No centro de momento, $\sqrt{s}$ é a energia total. Em
experimentos de altas energias, essa é a quantidade que determina quais
partículas podem ser produzidas.

\begin{activity}
Compare duas maneiras de obter uma grande energia de centro de massa: colidir
duas partículas de mesma energia em sentidos opostos ou lançar uma partícula
contra um alvo em repouso. Use conservação de quadrimomento para mostrar por
que colisores são muito mais eficientes que experimentos de alvo fixo.
\end{activity}

\section{Partículas massivas e partículas sem massa}

Para partículas massivas, existe referencial de repouso. Nesse referencial
$\bm p=0$ e $E=m$. Para partículas sem massa, não há referencial de repouso:
qualquer transformação de Lorentz admissível ainda observa a partícula se
movendo à velocidade da luz. Essa diferença é geométrica. O quadrimomento de
uma partícula massiva é temporal; o de um fóton é nulo.

Essa distinção reaparece em relatividade geral. Partículas massivas seguem
geodésicas temporais, parametrizadas naturalmente pelo tempo próprio. Fótons
seguem geodésicas nulas, para as quais o tempo próprio ao longo da trajetória é
zero. Isso não significa que ``a luz não exista durante o percurso'', mas sim
que não há relógio físico comóvel com um fóton.

\section{Problemas orientados}

\begin{enumerate}
  \item \textbf{Normalização da quadrivelocidade.} Parta de
  $U^\mu=\gamma(1,\bm v)$ e mostre que $U^\mu U_\mu=-1$. Depois derive
  $A^\mu U_\mu=0$ diferenciando essa normalização em relação a $\tau$.

  \item \textbf{Energia de limiar.} Uma partícula de massa $m$ atinge uma
  partícula idêntica em repouso e produz uma partícula de massa $M$ além dos
  produtos originais. Use o invariante $s$ para encontrar a energia mínima no
  laboratório.

  \item \textbf{Decaimento em dois corpos.} Uma partícula de massa $M$ decai em
  partículas de massas $m_1$ e $m_2$. No referencial de repouso inicial, derive
  as energias finais usando conservação de quadrimomento.

  \item \textbf{Sistema de fótons.} Dois fótons de energias $E_1$ e $E_2$ fazem
  ângulo $\theta$. Mostre que a massa invariante do sistema satisfaz
  \[
    M^2=2E_1E_2(1-\cos\theta).
  \]
  Discuta os casos $\theta=0$ e $\theta=\pi$.

  \item \textbf{Doppler covariante.} Use $E=-p_\mu U^\mu$ para obter a mudança
  de frequência observada quando fonte e observador se movem ao longo da mesma
  linha. Compare com a fórmula clássica no limite $v\ll1$.

  \item \textbf{Colisão elástica.} No centro de momento, duas partículas
  idênticas colidem elasticamente. Explique por que as energias finais são
  iguais às iniciais e por que apenas as direções dos momentos mudam.
\end{enumerate}

\clearpage
\section{Exercícios}

\begin{exercise}
Mostre que a quadriaceleração $A^\mu=\dd U^\mu/\dd\tau$ satisfaz
$A^\mu U_\mu=0$.
\end{exercise}

\begin{exercise}
Um fóton de frequência $\nu$ é observado por dois observadores inerciais em
movimento relativo. Use $E=-p_\mu U^\mu$ para obter o efeito Doppler
relativístico.
\end{exercise}

\chapter{Tensores no espaço-tempo plano}

\begin{objetivos}
Dominar a notação de índices, distinguir vetores de um-formas e compreender a
métrica como o objeto que transforma uma coisa na outra.
\end{objetivos}

\section{Vetores, um-formas e contrações}

Um vetor $V$ age como derivada direcional; uma um-forma $\omega$ age sobre
vetores e retorna escalares. Em componentes,
\[
  \omega(V)=\omega_\mu V^\mu.
\]
Um tensor de tipo $(r,s)$ é multilinear: recebe $s$ vetores e $r$ um-formas, e
retorna um número. Em física, componentes tensoriais são úteis porque suas
transformações preservam a forma das equações.

Para uma caixa de ferramentas prática sobre manipulação de índices, deltas,
símbolos antissimétricos e contrações, consulte o Apêndice~\ref{ap:indices}.
Ele foi escrito para ser usado como apoio de cálculo enquanto este capítulo é
estudado.

\section{Métrica}

No espaço-tempo de Minkowski,
\[
  \eta_{\mu\nu}=\mathrm{diag}(-1,1,1,1).
\]
A métrica define o produto escalar
\[
  A\cdot B=\eta_{\mu\nu}A^\mu B^\nu
\]
e permite baixar índices:
\[
  A_\mu=\eta_{\mu\nu}A^\nu.
\]
Com a convenção $(-,+,+,+)$, se $A^\mu=(A^0,A^1,A^2,A^3)$, então
$A_\mu=(-A^0,A^1,A^2,A^3)$.

\section{Regra prática de índices}

Índices repetidos uma vez acima e uma vez abaixo são somados. Índices livres
devem aparecer de modo idêntico nos dois lados da equação. Por exemplo,
\[
  A^\mu=B^\mu+C^\mu
\]
é uma equação vetorial válida, enquanto $A^\mu=B^\nu$ é mal escrita, pois
$\mu$ e $\nu$ são índices livres diferentes.

\begin{activity}
Classifique cada expressão como escalar, vetor, um-forma, tensor de ordem dois
ou expressão inválida:
\[
  A_\mu B^\mu,\qquad
  T^{\mu\nu}V_\nu,\qquad
  R_{\mu\nu}A^\rho,\qquad
  A_\mu+B^\mu.
\]
\end{activity}

\section{O que significa uma equação tensorial}

Uma equação tensorial bem escrita é uma afirmação independente de coordenadas.
Isso não quer dizer que as componentes sejam iguais em todos os referenciais.
Quer dizer que, se a relação entre objetos geométricos é verdadeira, qualquer
observador que escreva a relação em suas próprias componentes obterá uma
equação da mesma forma. Essa é a razão profunda para preferirmos tensores em
relatividade: eles separam física de descrição.

Por exemplo, a equação
\[
  p^\mu p_\mu=-m^2
\]
é escalar. Cada observador pode atribuir valores diferentes a $E$ e $\bm p$,
mas todos obtêm o mesmo valor para $E^2-\bm p^2$. Já a equação
\[
  A^\mu=B^\mu+C^\mu
\]
é vetorial: as componentes mudam de referencial, mas a relação entre os
vetores permanece. Em contraste, uma expressão como $A^\mu=B_\mu$ não é uma
igualdade tensorial natural sem uma métrica especificada, pois os dois lados
têm tipos diferentes.

\section{Mudança de base e transformação de componentes}

Considere uma mudança linear de coordenadas
\[
  x^{\mu'}=\Lambda^{\mu'}{}_\nu x^\nu.
\]
As componentes de um vetor transformam como
\[
  V^{\mu'}=\Lambda^{\mu'}{}_\nu V^\nu.
\]
As componentes de uma um-forma transformam com a matriz inversa:
\[
  \omega_{\mu'}=\frac{\partial x^\nu}{\partial x^{\mu'}}\omega_\nu.
\]
Essa diferença é exatamente o que torna a contração $\omega_\mu V^\mu$
invariante. A contração é a operação que junta um índice covariante e um
contravariante, eliminando a dependência da base.

\begin{example}
Em duas dimensões euclidianas, seja $\omega=(2,1)$ e $V=(3,4)$ em uma base
cartesiana. A contração é $\omega(V)=2\cdot3+1\cdot4=10$. Se mudarmos para uma
base rotacionada, as componentes individuais mudam, mas o número obtido pela
ação da um-forma sobre o vetor permanece o mesmo.
\end{example}

\section{A métrica como instrumento de medida}

A métrica tem duas funções simultâneas. Primeiro, ela mede produtos escalares e
normas. Segundo, ela fornece uma identificação entre vetores e um-formas. Essa
identificação não é automática em uma variedade; ela é informação geométrica
adicional. Em relatividade geral, quando a métrica passa a depender da posição,
essa informação geométrica torna-se dinâmica.

Em notação de índices, baixar um índice parece uma operação simples:
\[
  V_\mu=g_{\mu\nu}V^\nu.
\]
Mas o significado é físico. O objeto $V_\mu$ é a um-forma que, ao agir sobre
um vetor $W^\mu$, retorna o produto escalar $V\cdot W$. Essa leitura evita a
ideia equivocada de que índices altos e baixos são meras posições tipográficas.

\section{Erros comuns com índices}

Há três erros frequentes. O primeiro é repetir um índice mais de duas vezes em
um mesmo termo, como em $A_\mu B^\mu C^\mu$. O segundo é mudar o nome de um
índice livre entre os lados de uma equação. O terceiro é somar objetos de tipos
diferentes, como $A_\mu+B^\mu$, sem uma métrica ou identificação explícita.

Um bom hábito é ler os índices em voz alta: índices repetidos desaparecem por
soma; índices livres nomeiam o tipo do objeto resultante. Se depois de uma
contração ainda sobram índices $\alpha$ e $\beta$, o resultado é um tensor com
esses índices livres. Essa regra simples previne a maioria dos erros técnicos
no início do curso.

\section{Problemas orientados}

\begin{enumerate}
  \item \textbf{Contrações possíveis.} Dado $T^{\mu\nu}$, $A_\mu$ e $B^\mu$,
  liste todas as contrações escalares que podem ser formadas com um único
  $T$, um $A$ e um $B$. Explique quais expressões deixam índices livres.

  \item \textbf{Parte simétrica e antissimétrica.} Para um tensor $M_{\mu\nu}$,
  defina
  \[
    M_{(\mu\nu)}=\frac12(M_{\mu\nu}+M_{\nu\mu}),\qquad
    M_{[\mu\nu]}=\frac12(M_{\mu\nu}-M_{\nu\mu}).
  \]
  Mostre que $M_{\mu\nu}=M_{(\mu\nu)}+M_{[\mu\nu]}$.

  \item \textbf{Traço invariante.} Use a lei de transformação de um tensor
  misto $T^\mu{}_\nu$ para mostrar que $T^\mu{}_\mu$ não depende da base.

  \item \textbf{Gradiente.} Seja $f=t^2-x^2$. Calcule $\partial_\mu f$ e
  depois eleve o índice com a métrica de Minkowski. Interprete por que o sinal
  da componente temporal muda.

  \item \textbf{Produto externo.} Dados dois vetores $A^\mu$ e $B^\nu$,
  construa $A^\mu B^\nu$. Separe sua parte simétrica e antissimétrica.

  \item \textbf{Checagem de equações.} Verifique se as expressões
  $A_\mu=B_\mu+C_\mu$, $A^\mu B_\mu=C^\nu D_\nu$ e
  $T^{\mu\nu}=A^\mu B^\nu$ são tensorialmente consistentes.
\end{enumerate}

\section{Exemplos resolvidos adicionais}

\begin{example}
\textbf{Abaixamento de índices.} Seja $V^\mu=(3,2,0,0)$ em Minkowski com
assinatura $(-,+,+,+)$. Então
\[
  V_\mu=\eta_{\mu\nu}V^\nu=(-3,2,0,0).
\]
A norma é
\[
  V^\mu V_\mu=-9+4=-5.
\]
O vetor é temporal. Se fosse quadrimomento, corresponderia a uma partícula
massiva com massa invariante $\sqrt{5}$.
\end{example}

\begin{example}
\textbf{Contração de tensor antissimétrico.} Se $F_{\mu\nu}=-F_{\nu\mu}$ e
$S^{\mu\nu}=S^{\nu\mu}$, então
\[
  F_{\mu\nu}S^{\mu\nu}
  =F_{\nu\mu}S^{\nu\mu}
  =-F_{\mu\nu}S^{\nu\mu}
  =-F_{\mu\nu}S^{\mu\nu}.
\]
Logo a contração é igual ao seu negativo e, portanto, vale zero.
\end{example}

\begin{example}
\textbf{Transformação de uma um-forma.} Em duas dimensões, tome
$x'=2x$, $y'=y$. Se $\omega=\dd x$, então
\[
  \dd x=\frac{1}{2}\dd x'.
\]
Assim, a componente da um-forma na coordenada $x'$ é metade da componente
original. Vetores e um-formas transformam de modos inversos para preservar
$\omega(V)$.
\end{example}

\begin{example}
\textbf{Tensor como máquina multilinear.} Um tensor $T_{\mu\nu}$ recebe dois
vetores e devolve um número:
\[
  T(A,B)=T_{\mu\nu}A^\mu B^\nu.
\]
Se fixamos apenas o segundo vetor, obtemos uma um-forma:
\[
  \alpha_\mu=T_{\mu\nu}B^\nu.
\]
Essa leitura ajuda a lembrar que contração parcial reduz o tipo tensorial do
objeto.
\end{example}

\clearpage
\section{Exercícios}

\begin{exercise}
Mostre que o traço $T^\mu{}_\mu$ é invariante sob mudança de base.
\end{exercise}

\begin{exercise}
Se $F_{\mu\nu}=-F_{\nu\mu}$ e $S^{\mu\nu}=S^{\nu\mu}$, mostre que
$F_{\mu\nu}S^{\mu\nu}=0$.
\end{exercise}

\chapter{Fluidos relativísticos}

\begin{objetivos}
Introduzir o tensor energia-momento como fonte da gravidade e como objeto que
resume densidade de energia, fluxo de energia, densidade de momento e tensões.
\end{objetivos}

\section{Poeira}

Uma poeira é um conjunto de partículas sem pressão em seu referencial local de
repouso. Se $\rho$ é a densidade de energia própria e $U^\mu$ a
quadrivelocidade, então
\[
  T^{\mu\nu}=\rho U^\mu U^\nu.
\]
No referencial comóvel, $U^\mu=(1,0,0,0)$ e apenas $T^{00}=\rho$ é diferente
de zero.

\section{Fluido perfeito}

Um fluido perfeito isotrópico tem densidade de energia $\rho$ e pressão $p$.
Seu tensor energia-momento é
\[
  T^{\mu\nu}=(\rho+p)U^\mu U^\nu+p g^{\mu\nu}.
\]
No espaço plano, basta tomar $g^{\mu\nu}=\eta^{\mu\nu}$. Em relatividade geral,
o mesmo formato vale localmente em um espaço-tempo curvo.

\begin{resultado}
A conservação local de energia e momento é
\[
  \nabla_\mu T^{\mu\nu}=0.
\]
Em espaço-tempo plano e coordenadas cartesianas, isso se reduz a
$\partial_\mu T^{\mu\nu}=0$.
\end{resultado}

\section{Como ler o tensor energia-momento}

O tensor energia-momento é uma tabela de fluxos. A componente $T^{00}$ é
densidade de energia. As componentes $T^{0i}$ representam fluxo de energia na
direção $i$. As componentes $T^{i0}$ representam densidade de momento. As
componentes $T^{ij}$ representam fluxo do momento $i$ através de uma superfície
normal à direção $j$; por isso incluem pressão e tensões.

No caso de um fluido perfeito em repouso local,
\[
  T^{\mu\nu}=
  \begin{pmatrix}
  \rho&0&0&0\\
  0&p&0&0\\
  0&0&p&0\\
  0&0&0&p
  \end{pmatrix}.
\]
Essa matriz mostra por que pressão gravita em relatividade geral. A fonte da
geometria não é apenas massa no sentido newtoniano, mas todo o tensor
energia-momento. Em situações não relativísticas, $p\ll\rho$ e a pressão pode
ser desprezada como fonte gravitacional. Em estrelas de nêutrons, radiação no
Universo primordial ou campos relativísticos, essa aproximação deixa de ser
segura.

\section{Equações de estado}

Para fechar um problema de fluido, é necessário relacionar pressão e densidade.
Uma forma simples é
\[
  p=w\rho.
\]
Poeira fria tem $w=0$. Radiação tem $w=1/3$, pois o traço do tensor
energia-momento eletromagnético é nulo. Uma constante cosmológica efetiva tem
$w=-1$. Esses valores reaparecerão em cosmologia, onde determinam como a
densidade de energia evolui com a expansão.

\begin{example}
No referencial comóvel de um fluido de radiação,
\[
  T^\mu{}_\mu=-\rho+3p.
\]
Para radiação, $T^\mu{}_\mu=0$, logo $p=\rho/3$. Essa é uma maneira rápida de
lembrar a equação de estado relativística para um gás de partículas sem massa.
\end{example}

\section{Conservação local e leis integrais}

Em espaço plano, a equação $\partial_\mu T^{\mu\nu}=0$ pode ser integrada em
um volume. A componente $\nu=0$ expressa conservação de energia: a variação da
energia dentro do volume é igual ao fluxo líquido através da fronteira. As
componentes espaciais expressam conservação de momento. Em espaço-tempo curvo,
a forma local torna-se $\nabla_\mu T^{\mu\nu}=0$, mas a passagem para uma
energia global conservada exige simetria temporal. Essa sutileza é uma das
diferenças profundas entre gravitação relativística e teorias de campo em um
fundo fixo.

\begin{activity}
Escolha três sistemas físicos: poeira fria, gás relativístico e energia escura.
Para cada um, escreva uma matriz diagonal plausível para $T^{\mu}{}_{\nu}$ no
referencial comóvel e discuta o sinal da pressão.
\end{activity}

\section{Problemas orientados}

\begin{enumerate}
  \item \textbf{Poeira em movimento.} Para uma poeira com densidade própria
  $\rho$ e velocidade $v$ na direção $x$, calcule explicitamente os componentes
  $T^{00}$, $T^{0x}$ e $T^{xx}$ em Minkowski.

  \item \textbf{Traço do fluido perfeito.} Mostre que
  \[
    T^\mu{}_\mu=-\rho+3p.
  \]
  Aplique o resultado a poeira, radiação e constante cosmológica.

  \item \textbf{Limite não relativístico.} Assuma $p\ll\rho$ e $v\ll1$.
  Identifique os termos dominantes de $T^{\mu\nu}$ e explique por que a fonte
  newtoniana é aproximadamente a densidade de massa.

  \item \textbf{Conservação de partículas.} Partindo de $N^\mu=nU^\mu$,
  escreva $\partial_\mu N^\mu=0$ em um referencial onde o fluido tem velocidade
  ordinária $\bm v$.

  \item \textbf{Pressão como fluxo.} Explique por que $T^{xx}=p$ no referencial
  comóvel significa fluxo de momento $x$ através de uma superfície normal a
  $x$. Use colisões microscópicas com uma parede como imagem física.

  \item \textbf{Som de um fluido.} Para uma equação de estado $p=p(\rho)$, a
  velocidade do som relativística satisfaz $c_s^2=\dd p/\dd\rho$. Discuta por
  que causalidade sugere $c_s\le1$.
\end{enumerate}

\clearpage
\section{Exercícios}

\begin{exercise}
Calcule $T^{\mu\nu}$ de um fluido perfeito no referencial de repouso local.
\end{exercise}

\begin{exercise}
No limite não relativístico, identifique quais componentes de $T^{\mu\nu}$ são
dominantes para uma poeira lenta.
\end{exercise}

\chapter{Da gravidade à curvatura}

\begin{objetivos}
Preparar a transição conceitual: gravidade não é uma força em um palco fixo,
mas manifestação da geometria do espaço-tempo.
\end{objetivos}

\section{O princípio da equivalência}

Localmente, um campo gravitacional uniforme pode ser eliminado pela escolha de
um referencial em queda livre. Essa afirmação não diz que a gravidade desaparece
globalmente. Gradientes do campo gravitacional, isto é, efeitos de maré, não
podem ser removidos em uma região finita.

\begin{resultado}
O conteúdo físico da curvatura está nos efeitos de maré. A conexão pode ser
zerada em um ponto por uma escolha de coordenadas; o tensor de curvatura não.
\end{resultado}

\section{Coordenadas curvilíneas}

Mesmo no plano euclidiano, coordenadas polares produzem símbolos de Christoffel
não nulos. Com
\[
  \dd \ell^2=\dd r^2+r^2\dd\phi^2,
\]
os símbolos não nulos são
\[
  \Gamma^r_{\phi\phi}=-r,\qquad
  \Gamma^\phi_{r\phi}=\Gamma^\phi_{\phi r}=\frac{1}{r}.
\]
Isso ensina uma lição importante: Christoffel não nulo não significa curvatura
não nula. Ele pode ser apenas efeito de coordenadas.

\section{Equação geodésica}

Uma geodésica satisfaz
\[
  \frac{\dd^2 x^\mu}{\dd\lambda^2}
  +\Gamma^\mu_{\alpha\beta}
  \frac{\dd x^\alpha}{\dd\lambda}
  \frac{\dd x^\beta}{\dd\lambda}=0,
\]
onde $\lambda$ é parâmetro afim. Para partículas massivas, podemos escolher
$\lambda=\tau$. Para luz, usa-se um parâmetro afim nulo.

\section{Christoffel não é força}

É tentador interpretar $\Gamma^\mu_{\alpha\beta}$ como força gravitacional,
pois ele aparece na equação geodésica em uma posição parecida com aceleração.
Essa interpretação é perigosa. Os símbolos de Christoffel não formam um tensor
e podem ser anulados em um ponto por escolha adequada de coordenadas. Uma
quantidade física local não deve depender desse tipo de escolha.

A maneira correta de ler a equação geodésica é: uma partícula livre mantém seu
vetor tangente paralelamente transportado ao longo da própria trajetória. Em
coordenadas curvilíneas, a mudança dos vetores de base introduz termos
adicionais. Em uma geometria curva, termos semelhantes descrevem também a
estrutura gravitacional local. A distinção entre efeito de coordenadas e
efeito físico é feita pela curvatura.

\section{Derivação por princípio variacional}

Para uma curva $x^\mu(\lambda)$, considere a Lagrangiana
\[
  L=\frac{1}{2}g_{\mu\nu}(x)\dot x^\mu\dot x^\nu,
\]
onde ponto denota derivada em relação a $\lambda$. As equações de
Euler-Lagrange são
\[
  \frac{\dd}{\dd\lambda}\left(\frac{\partial L}{\partial \dot x^\rho}\right)
  -\frac{\partial L}{\partial x^\rho}=0.
\]
Como
\[
  \frac{\partial L}{\partial \dot x^\rho}=g_{\rho\nu}\dot x^\nu,
\]
temos
\[
  g_{\rho\nu}\ddot x^\nu+\partial_\alpha g_{\rho\nu}\dot x^\alpha\dot x^\nu
  -\frac{1}{2}\partial_\rho g_{\mu\nu}\dot x^\mu\dot x^\nu=0.
\]
Multiplicando por $g^{\rho\sigma}$ e rearranjando índices, obtemos
\[
  \ddot x^\sigma+\Gamma^\sigma_{\mu\nu}\dot x^\mu\dot x^\nu=0.
\]
Assim, geodésicas podem ser vistas tanto como curvas de transporte paralelo do
tangente quanto como extremantes do comprimento ou do tempo próprio.

\section{Marés como sinal de curvatura}

Imagine duas partículas em queda livre próximas uma da outra. Em um campo
gravitacional uniforme ideal, elas mantêm separação constante em um referencial
em queda livre. No campo real da Terra, duas partículas separadas radialmente
tendem a se afastar, enquanto partículas separadas transversalmente tendem a se
aproximar. Esse padrão é chamado efeito de maré.

Na linguagem relativística, marés são descritas por curvatura. A conexão pode
ser removida em um ponto; diferenças da conexão em pontos vizinhos não podem
ser removidas em uma região finita. Por isso a curvatura, e não a conexão,
contém a parte observável da gravidade local.

\begin{example}
Em coordenadas polares no plano, os símbolos de Christoffel são não nulos, mas
se transportarmos um vetor ao longo de um pequeno circuito fechado e retornarmos
ao ponto inicial, o vetor não sofre rotação líquida associada à curvatura do
plano. Em uma esfera, o mesmo experimento produz uma diferença angular
proporcional à área envolvida. Essa é uma manifestação geométrica da curvatura.
\end{example}

\section{Problemas orientados}

\begin{enumerate}
  \item \textbf{Inversa da métrica polar.} Para
  $\dd\ell^2=\dd r^2+r^2\dd\phi^2$, escreva $g_{ij}$ e $g^{ij}$. Use a fórmula
  de Christoffel para obter os símbolos não nulos.

  \item \textbf{Geodésica radial.} Mostre que curvas com $\phi=\mathrm{const}$
  são geodésicas no plano polar. Interprete o resultado geometricamente.

  \item \textbf{Círculos não são geodésicas.} Tome $r=R$ constante e
  $\phi=\lambda$. Substitua na equação geodésica e mostre que a equação radial
  não é satisfeita.

  \item \textbf{Esfera unitária.} Para
  $\dd\ell^2=\dd\theta^2+\sin^2\theta\,\dd\phi^2$, calcule
  $\Gamma^\theta_{\phi\phi}$ e $\Gamma^\phi_{\theta\phi}$.

  \item \textbf{Grande círculos.} Use simetria para argumentar que o equador de
  uma esfera é geodésica. Depois verifique pela equação geodésica.

  \item \textbf{Maré newtoniana.} Para $\Phi=-GM/r$, calcule qualitativamente
  as segundas derivadas radiais e transversais. Relacione-as a alongamento
  radial e compressão transversal.
\end{enumerate}

\clearpage
\section{Exercícios}

\begin{exercise}
Use os símbolos de Christoffel polares para mostrar que uma reta no plano
continua sendo geodésica, embora suas coordenadas polares satisfaçam equações
não triviais.
\end{exercise}

\begin{exercise}
Explique por que uma estação espacial em órbita está em queda livre, apesar de
sua distância ao centro da Terra ser aproximadamente constante.
\end{exercise}

\chapter{Variedades, derivada covariante e curvatura}

\begin{objetivos}
Formalizar o cálculo tensorial em espaços-tempos curvos e obter os objetos
geométricos que entram nas equações de Einstein.
\end{objetivos}

\section{Variedades e métricas}

Uma variedade diferenciável é um espaço que, localmente, parece $\mathbb{R}^n$.
O espaço-tempo da relatividade geral é uma variedade quadridimensional equipada
com uma métrica lorentziana $g_{\mu\nu}$.

O elemento de linha é
\[
  \dd s^2=g_{\mu\nu}\dd x^\mu\dd x^\nu.
\]
Quando $g_{\mu\nu}$ depende da posição, a geometria muda de ponto a ponto.

\section{Derivada covariante}

Para um vetor,
\[
  \nabla_\mu V^\nu=\partial_\mu V^\nu+\Gamma^\nu_{\mu\alpha}V^\alpha.
\]
Para uma um-forma,
\[
  \nabla_\mu \omega_\nu=\partial_\mu \omega_\nu-\Gamma^\alpha_{\mu\nu}\omega_\alpha.
\]
O sinal muda porque a contração $\omega_\nu V^\nu$ deve ter derivada usual
compatível.

Para a conexão de Levi-Civita, impomos torção nula e compatibilidade métrica:
\[
  \Gamma^\rho_{\mu\nu}=\Gamma^\rho_{\nu\mu},
  \qquad
  \nabla_\alpha g_{\mu\nu}=0.
\]
Daí segue
\[
  \Gamma^\rho_{\mu\nu}
  =\frac{1}{2}g^{\rho\sigma}
  \left(\partial_\mu g_{\nu\sigma}
  +\partial_\nu g_{\mu\sigma}
  -\partial_\sigma g_{\mu\nu}\right).
\]

\section{Curvatura}

O tensor de Riemann mede a falha de derivadas covariantes em comutar:
\[
  [\nabla_\mu,\nabla_\nu]V^\rho
  =R^\rho{}_{\sigma\mu\nu}V^\sigma.
\]
Em componentes,
\[
  R^\rho{}_{\sigma\mu\nu}
  =\partial_\mu\Gamma^\rho_{\nu\sigma}
  -\partial_\nu\Gamma^\rho_{\mu\sigma}
  +\Gamma^\rho_{\mu\lambda}\Gamma^\lambda_{\nu\sigma}
  -\Gamma^\rho_{\nu\lambda}\Gamma^\lambda_{\mu\sigma}.
\]
As contrações relevantes são
\[
  R_{\mu\nu}=R^\rho{}_{\mu\rho\nu},\qquad
  R=g^{\mu\nu}R_{\mu\nu},
\]
e o tensor de Einstein é
\[
  G_{\mu\nu}=R_{\mu\nu}-\frac{1}{2}Rg_{\mu\nu}.
\]

\begin{resultado}
A identidade de Bianchi contraída implica
\[
  \nabla_\mu G^{\mu\nu}=0.
\]
Essa propriedade geométrica antecipa a conservação local de energia e momento
da matéria.
\end{resultado}

\section{Compatibilidade métrica}

A condição $\nabla_\alpha g_{\mu\nu}=0$ afirma que o transporte paralelo
preserva produtos escalares. Se dois vetores são transportados paralelamente ao
longo de uma curva, o produto escalar entre eles permanece constante. Essa é
uma exigência natural para uma conexão que deve representar a geometria
determinada pela métrica.

Junto com torção nula, a compatibilidade métrica determina de modo único a
conexão de Levi-Civita. Em relatividade geral padrão, essa é a conexão usada
para definir geodésicas, curvatura e derivada covariante. Teorias mais gerais
podem relaxar torção nula ou compatibilidade métrica, mas isso está além do
curso introdutório.

\section{Simetrias do tensor de Riemann}

O tensor de Riemann possui várias simetrias. Com todos os índices abaixados,
\[
  R_{\rho\sigma\mu\nu}
  =-R_{\sigma\rho\mu\nu}
  =-R_{\rho\sigma\nu\mu}
  =R_{\mu\nu\rho\sigma}.
\]
Além disso, vale a primeira identidade de Bianchi,
\[
  R_{\rho[\sigma\mu\nu]}=0.
\]
Essas simetrias reduzem drasticamente o número de componentes independentes.
Em quatro dimensões, o tensor de Riemann possui 20 componentes independentes,
não 256. O tensor de Ricci contém 10 componentes independentes. As 10 restantes
estão associadas ao tensor de Weyl, que descreve a parte livre da curvatura,
incluindo marés de vácuo e ondas gravitacionais.

\section{Curvatura intrínseca e extrínseca}

Uma superfície pode parecer curva por estar embutida em um espaço maior, mas a
curvatura relevante em relatividade geral é intrínseca. Habitantes
bidimensionais de uma esfera poderiam detectar sua curvatura medindo somas de
ângulos de triângulos, sem olhar para uma terceira dimensão. Do mesmo modo, a
curvatura do espaço-tempo não exige imaginá-lo embutido em um espaço externo.
Ela é medida por experimentos internos: transporte paralelo, desvio geodésico,
volumes de pequenas bolas geodésicas e propagação de luz.

\begin{example}
Em uma esfera de raio $a$, um triângulo formado por dois meridianos separados
por $90^\circ$ e pelo equador possui três ângulos retos. A soma é $270^\circ$,
maior que $180^\circ$. Essa diferença é uma medida integrada da curvatura
positiva da esfera.
\end{example}

\section{Desvio geodésico}

Seja $\xi^\mu$ o vetor separação entre duas geodésicas vizinhas e $U^\mu$ o
vetor tangente à família de geodésicas. A equação de desvio geodésico é
\[
  \frac{\D^2\xi^\mu}{\D\tau^2}
  =
  -R^\mu{}_{\alpha\nu\beta}U^\alpha\xi^\nu U^\beta.
\]
Essa equação é uma das formas mais diretas de interpretar curvatura. Ela diz
que a aceleração relativa entre partículas em queda livre é controlada pelo
tensor de Riemann. No limite newtoniano, ela se reduz à equação de maré
envolvendo segundas derivadas do potencial gravitacional.

\begin{activity}
Considere duas partículas inicialmente em repouso relativo, separadas por uma
pequena distância radial no campo de uma massa central. Use intuição
newtoniana para determinar se a separação aumenta ou diminui. Depois associe o
resultado ao sinal de componentes apropriadas de curvatura.
\end{activity}

\section{Problemas orientados}

\begin{enumerate}
  \item \textbf{Compatibilidade métrica.} Partindo de $\nabla_\alpha g_{\mu\nu}=0$,
  mostre que $\nabla_\alpha g^{\mu\nu}=0$. A pista é derivar
  $g_{\mu\rho}g^{\rho\nu}=\delta_\mu{}^\nu$.

  \item \textbf{Divergência de vetor.} Mostre que
  \[
    \nabla_\mu V^\mu=\frac{1}{\sqrt{|g|}}\partial_\mu(\sqrt{|g|}V^\mu).
  \]
  Esse resultado é útil para leis de conservação em coordenadas gerais.

  \item \textbf{Comutador em escalares.} Verifique que
  $[\nabla_\mu,\nabla_\nu]f=0$ para qualquer escalar $f$, assumindo conexão sem
  torção.

  \item \textbf{Ricci da esfera.} Usando os Christoffel da esfera unitária,
  calcule pelo menos uma componente de $R_{\mu\nu}$ e compare com
  $R_{\mu\nu}=g_{\mu\nu}$ para esfera unitária.

  \item \textbf{Número de componentes.} Conte os componentes independentes de
  um tensor simétrico $4\times4$. Depois discuta por que as equações de
  Einstein têm dez componentes, mas nem todas são evolutivas independentes.

  \item \textbf{Curvatura zero.} Explique por que todos os componentes do
  tensor de Riemann se anulam em Minkowski, mesmo quando usamos coordenadas
  polares ou esféricas.
\end{enumerate}

\section{Exemplos resolvidos adicionais}

\begin{example}
\textbf{Divergência em coordenadas polares.} No plano,
$g_{ij}=\mathrm{diag}(1,r^2)$ e $\sqrt{|g|}=r$. Para um vetor
$V^i=(V^r,V^\phi)$,
\[
  \nabla_iV^i=\frac{1}{r}\partial_r(rV^r)+\frac{1}{r}\partial_\phi(rV^\phi).
\]
Como $r$ não depende de $\phi$,
\[
  \nabla_iV^i=\partial_rV^r+\frac{1}{r}V^r+\partial_\phi V^\phi.
\]
Esse é o termo extra familiar da divergência em polares, agora entendido como
consequência de $\sqrt{|g|}$.
\end{example}

\begin{example}
\textbf{Geodésicas na esfera.} Na esfera unitária,
\[
  \dd\ell^2=\dd\theta^2+\sin^2\theta\,\dd\phi^2.
\]
Os símbolos relevantes são
\[
  \Gamma^\theta_{\phi\phi}=-\sin\theta\cos\theta,\qquad
  \Gamma^\phi_{\theta\phi}=\cot\theta.
\]
No equador, $\theta=\pi/2$ e $\dot\theta=0$. A equação para $\theta$ contém
$-\sin\theta\cos\theta\,\dot\phi^2$, que se anula no equador. Portanto o
equador é geodésica.
\end{example}

\begin{example}
\textbf{Curvatura por transporte paralelo.} Transporte um vetor ao longo de um
triângulo pequeno em uma superfície curva. A diferença entre o vetor inicial e
final é proporcional à área do circuito e ao tensor de curvatura. No limite de
área infinitesimal, essa relação fornece uma interpretação geométrica direta de
$R^\rho{}_{\sigma\mu\nu}$: ele mede a rotação infinitesimal produzida por um
paralelogramo infinitesimal.
\end{example}

\begin{example}
\textbf{Coordenadas normais.} Em qualquer ponto de uma variedade pseudo-
riemanniana, é possível escolher coordenadas tais que
$g_{\mu\nu}=\eta_{\mu\nu}$ e $\partial_\alpha g_{\mu\nu}=0$ naquele ponto. Como
os Christoffel dependem de primeiras derivadas da métrica, eles se anulam ali.
As segundas derivadas, porém, permanecem e codificam a curvatura.
\end{example}

\clearpage
\section{Exercícios}

\begin{exercise}
Calcule os símbolos de Christoffel da métrica da esfera unitária,
$\dd \ell^2=\dd\theta^2+\sin^2\theta\,\dd\phi^2$.
\end{exercise}

\begin{exercise}
Mostre que, em coordenadas localmente inerciais em um ponto, podemos ter
$\Gamma^\rho_{\mu\nu}=0$ nesse ponto, mas em geral não podemos anular
$\partial_\alpha\Gamma^\rho_{\mu\nu}$.
\end{exercise}

\chapter{Física em espaço-tempo curvo}

\begin{objetivos}
Aprender como partículas, luz, relógios e observadores são descritos em uma
geometria curva.
\end{objetivos}

\section{Movimento livre}

Na relatividade geral, uma partícula livre segue uma geodésica temporal:
\[
  U^\nu\nabla_\nu U^\mu=0.
\]
Essa equação expressa que o vetor tangente é transportado paralelamente ao
longo da própria curva.

Para luz,
\[
  k^\nu\nabla_\nu k^\mu=0,
  \qquad
  k^\mu k_\mu=0.
\]

\section{Energia medida por um observador}

Se uma partícula tem quadrimomento $p^\mu$ e o observador tem
quadrivelocidade $U^\mu$, a energia medida localmente é
\[
  E_{\rm obs}=-p_\mu U^\mu.
\]
Essa fórmula é uma das mais úteis da relatividade: ela vale em espaço plano,
em espaço curvo, para partículas massivas e para fótons.

\section{Simetrias e quantidades conservadas}

Se $\xi^\mu$ é um vetor de Killing, então ao longo de uma geodésica com
momento $p^\mu$,
\[
  C=\xi_\mu p^\mu
\]
é constante. Em uma geometria estacionária, a simetria temporal fornece uma
energia conservada. Em uma geometria axial, a simetria angular fornece momento
angular conservado.

\section{Observadores não são coordenadas}

Em relatividade geral, coordenadas são rótulos; observadores são objetos
físicos. Um observador ideal é representado por uma linha de mundo temporal e
por uma quadrivelocidade $U^\mu$ normalizada. Dois observadores podem usar as
mesmas coordenadas e ainda assim medir energias diferentes, pois a medida
depende de sua quadrivelocidade. Do mesmo modo, uma coordenada temporal $t$ não
é automaticamente o tempo próprio de alguém.

Essa distinção é crucial em espaços-tempos curvos. Na métrica de
Schwarzschild, por exemplo, o observador estático em raio fixo mede um tempo
próprio relacionado a $\dd t$ por um fator que depende de $r$. Um observador em
queda livre atravessando o mesmo ponto possui outra quadrivelocidade e mede
outra decomposição local de energia e momento.

\section{Tetradas locais}

Uma tetrada é um conjunto de quatro vetores ortonormais carregados por um
observador. O vetor temporal da tetrada é a quadrivelocidade do observador; os
três vetores espaciais definem seus eixos locais. Em componentes de tetrada, a
métrica parece a métrica de Minkowski:
\[
  g(e_{\hat\alpha},e_{\hat\beta})=\eta_{\hat\alpha\hat\beta}.
\]
Índices com chapéu indicam componentes medidas localmente, não componentes em
uma carta qualquer.

Essa linguagem é útil porque aparelhos reais medem componentes locais. Um
detector não mede diretamente $p_t$ ou $p_r$ como números abstratos de uma
carta; ele mede energia e direção em seu laboratório. A conversão entre
componentes coordenadas e componentes medidas passa pela tetrada.

\section{Constantes de movimento e energia local}

Em uma geometria estacionária, $p_t$ pode ser constante ao longo de uma
geodésica. Isso não significa que todos os observadores locais meçam a mesma
energia. A energia conservada associada ao infinito e a energia local medida
por um observador são conceitos relacionados, mas distintos:
\[
  E_{\rm loc}=-p_\mu U^\mu.
\]
No caso de um fóton subindo em um campo gravitacional estático, a quantidade
conservada ligada à simetria temporal permanece fixa, enquanto a energia
medida por observadores estáticos muda com o potencial gravitacional. Essa é a
origem do redshift gravitacional.

\begin{example}
Em espaço plano, dois observadores inerciais com quadrivelocidades $U^\mu$ e
$V^\mu$ medem energias diferentes para a mesma partícula. A diferença é
puramente cinemática e corresponde ao efeito Doppler. Em espaço-tempo curvo,
além dessa diferença cinemática, pode haver variação associada à posição dos
observadores na geometria.
\end{example}

\section{Problemas orientados}

\begin{enumerate}
  \item \textbf{Energia local.} Em Minkowski, tome $p^\mu=m\gamma(1,v,0,0)$ e
  um observador em repouso $U^\mu=(1,0,0,0)$. Calcule $E=-p_\mu U^\mu$.

  \item \textbf{Observador em movimento.} Repita o cálculo anterior para um
  observador com velocidade $u$ na direção $x$. Mostre que aparece a energia
  transformada por Lorentz.

  \item \textbf{Killing temporal.} Em uma métrica independente de $t$, mostre
  que $\xi^\mu=(1,0,0,0)$ é candidato natural a vetor de Killing. Escreva a
  constante associada para uma geodésica.

  \item \textbf{Killing axial.} Em uma métrica independente de $\phi$, obtenha
  a constante associada a rotações. Interprete-a como momento angular.

  \item \textbf{Geodésica nula.} Explique por que não podemos usar tempo
  próprio para parametrizar a trajetória de um fóton. O que significa escolher
  um parâmetro afim?

  \item \textbf{Tetrada estática.} Para uma métrica diagonal estática, construa
  qualitativamente uma tetrada de observador em repouso nas coordenadas
  espaciais.
\end{enumerate}

\clearpage
\section{Exercícios}

\begin{exercise}
Mostre que $C=\xi_\mu p^\mu$ é conservado ao longo de uma geodésica quando
$\nabla_{(\mu}\xi_{\nu)}=0$.
\end{exercise}

\chapter{Equações de Einstein}

\begin{objetivos}
Compreender a forma, a motivação e os limites físicos das equações de campo.
\end{objetivos}

\section{A equação de campo}

As equações de Einstein são
\[
  G_{\mu\nu}+\Lambda g_{\mu\nu}=8\pi G T_{\mu\nu},
\]
ou, em unidades geométricas $G=c=1$,
\[
  G_{\mu\nu}+\Lambda g_{\mu\nu}=8\pi T_{\mu\nu}.
\]
O lado esquerdo descreve geometria; o lado direito descreve matéria e energia.
A constante cosmológica $\Lambda$ pode ser interpretada geometricamente ou como
componente efetiva do conteúdo energético do Universo.

\section{Limite newtoniano}

No limite de campo fraco e velocidades pequenas, escrevemos
\[
  g_{00}\approx -(1+2\Phi),
\]
onde $\Phi$ é o potencial gravitacional newtoniano. As equações de Einstein
reduzem-se à equação de Poisson
\[
  \nabla^2\Phi=4\pi G\rho.
\]
Esse limite fixa o fator $8\pi G$ na equação de campo.

\section{Conservação local}

Como $\nabla_\mu G^{\mu\nu}=0$, as equações de Einstein exigem
\[
  \nabla_\mu T^{\mu\nu}=0.
\]
Essa conservação é local. Em espaço-tempos genéricos, não existe uma energia
gravitacional global simples e tensorialmente definida.

\section{Por que o tensor de Einstein}

Queremos uma equação que relacione geometria e matéria. O lado da matéria deve
ser o tensor energia-momento $T_{\mu\nu}$, pois ele contém densidade de energia,
fluxos e tensões. Portanto o lado geométrico deve ser também um tensor
simétrico de tipo $(0,2)$. Além disso, deve ter divergência covariante nula,
para ser compatível com conservação local de energia e momento. O tensor de
Einstein satisfaz exatamente essa exigência:
\[
  \nabla_\mu G^{\mu\nu}=0.
\]

Outro requisito é que a equação contenha no máximo segundas derivadas da
métrica. Isso é importante para que a teoria tenha um problema de valor inicial
bem comportado, análogo ao de outras teorias clássicas. O tensor de Einstein é
construído a partir da curvatura, que envolve segundas derivadas da métrica e
termos quadráticos em primeiras derivadas.

\section{Traço das equações de campo}

Sem constante cosmológica, as equações são
\[
  R_{\mu\nu}-\frac{1}{2}Rg_{\mu\nu}=8\pi G T_{\mu\nu}.
\]
Contraindo com $g^{\mu\nu}$ em quatro dimensões,
\[
  R-\frac{1}{2}R\,g^{\mu\nu}g_{\mu\nu}=8\pi G T.
\]
Como $g^{\mu\nu}g_{\mu\nu}=4$, obtemos
\[
  R-2R=8\pi G T,
  \qquad
  R=-8\pi G T.
\]
Substituindo de volta,
\[
  R_{\mu\nu}=8\pi G\left(T_{\mu\nu}-\frac{1}{2}Tg_{\mu\nu}\right).
\]
Essa forma é muitas vezes útil em cálculos. No vácuo, $T_{\mu\nu}=0$, ela
implica $R_{\mu\nu}=0$, embora o tensor de Riemann completo possa ser não nulo.
Por isso há gravidade de maré e ondas gravitacionais mesmo em regiões sem
matéria.

\section{O limite fraco com mais detalhe}

No limite newtoniano, assumimos campo estático, velocidades pequenas e métrica
quase plana. Escrevemos
\[
  g_{00}=-(1+2\Phi),\qquad |\Phi|\ll1.
\]
Para uma partícula lenta, a equação geodésica espacial é dominada por termos
com duas derivadas temporais da coordenada:
\[
  \frac{\dd^2 x^i}{\dd t^2}\approx-\Gamma^i_{00}.
\]
Como
\[
  \Gamma^i_{00}\approx \partial_i\Phi,
\]
temos
\[
  \frac{\dd^2 x^i}{\dd t^2}\approx-\partial_i\Phi,
\]
que é a equação newtoniana. A componente $00$ das equações de Einstein fornece
a equação de Poisson,
\[
  \nabla^2\Phi=4\pi G\rho.
\]

\section{Constante cosmológica como fluido}

Movendo o termo $\Lambda g_{\mu\nu}$ para o lado direito, podemos definir
\[
  T^{(\Lambda)}_{\mu\nu}=-\frac{\Lambda}{8\pi G}g_{\mu\nu}.
\]
No referencial comóvel de um fluido perfeito, isso corresponde a
\[
  \rho_\Lambda=\frac{\Lambda}{8\pi G},
  \qquad
  p_\Lambda=-\rho_\Lambda.
\]
Essa pressão negativa gera aceleração da expansão cosmológica nas equações de
Friedmann. A interpretação microscópica da energia escura permanece uma das
questões abertas da física fundamental.

\section{Problemas orientados}

\begin{enumerate}
  \item \textbf{Traço de Einstein.} Contraia as equações de campo em quatro
  dimensões e obtenha $R=-8\pi G T$. Depois escreva a forma equivalente para
  $R_{\mu\nu}$.

  \item \textbf{Vácuo não é plano.} Explique por que $R_{\mu\nu}=0$ não implica
  $R^\rho{}_{\sigma\mu\nu}=0$. Use o exterior de uma estrela esférica como
  exemplo.

  \item \textbf{Unidades.} Restaure fatores de $c$ e verifique que
  $G_{\mu\nu}=8\pi G T_{\mu\nu}/c^4$ tem dimensões consistentes.

  \item \textbf{Geodésica newtoniana.} Complete o cálculo de
  $\Gamma^i_{00}$ para $g_{00}=-(1+2\Phi)$ e obtenha
  $\ddot x^i=-\partial_i\Phi$.

  \item \textbf{Pressão gravitante.} Use o fluido perfeito para discutir por
  que pressão contribui para a fonte gravitacional em regimes relativísticos.

  \item \textbf{Energia gravitacional.} Discuta por que não existe, em geral,
  um tensor local de energia do campo gravitacional análogo ao tensor
  energia-momento da matéria.
\end{enumerate}

\section{Exemplos resolvidos adicionais}

\begin{example}
\textbf{Equações no vácuo.} Se $T_{\mu\nu}=0$ e $\Lambda=0$, então
\[
  R_{\mu\nu}-\frac12Rg_{\mu\nu}=0.
\]
Contraindo, obtemos $R-2R=0$, logo $R=0$. Substituindo de volta,
$R_{\mu\nu}=0$. Isso não torna o espaço-tempo necessariamente plano, pois a
parte de Weyl da curvatura pode permanecer.
\end{example}

\begin{example}
\textbf{Constante cosmológica em vácuo.} Se $T_{\mu\nu}=0$ mas
$\Lambda\ne0$, as equações dão
\[
  R_{\mu\nu}-\frac12Rg_{\mu\nu}+\Lambda g_{\mu\nu}=0.
\]
Contraindo em quatro dimensões,
\[
  -R+4\Lambda=0,\qquad R=4\Lambda.
\]
Logo
\[
  R_{\mu\nu}=\Lambda g_{\mu\nu}.
\]
Espaços de de Sitter e anti-de Sitter são soluções maximamente simétricas
desse tipo, com sinal de $\Lambda$ determinando o caráter da curvatura.
\end{example}

\begin{example}
\textbf{Estimativa de campo fraco.} Para a Terra,
$GM/(Rc^2)\sim7\times10^{-10}$. Para o Sol,
$GM/(Rc^2)\sim2\times10^{-6}$. Para uma estrela de nêutrons, essa quantidade
pode ser da ordem de $0{,}1$. Isso explica por que relatividade geral é uma
correção pequena em laboratórios terrestres, mas central para objetos compactos.
\end{example}

\begin{example}
\textbf{Pressão e traço.} Para radiação, $T=0$. A forma traçada das equações
mostra que $R=-8\pi GT=0$, mas isso não implica $R_{\mu\nu}=0$, pois
$T_{\mu\nu}$ pode ser não nulo mesmo com traço zero. Esse caso é importante no
Universo primordial dominado por radiação.
\end{example}

\clearpage
\section{Exercícios}

\begin{exercise}
Explique por que a equação $R_{\mu\nu}=8\pi T_{\mu\nu}$ não é a melhor
candidata para a equação de campo relativística.
\end{exercise}

\begin{exercise}
No vácuo com $\Lambda=0$, mostre que as equações de Einstein implicam
$R_{\mu\nu}=0$.
\end{exercise}

\chapter{Ondas gravitacionais}

\begin{objetivos}
Descrever ondas gravitacionais como perturbações da métrica, entender suas
polarizações e conectar a teoria linearizada à detecção interferométrica.
\end{objetivos}

\section{Campo fraco}

Em uma região quase plana,
\[
  g_{\mu\nu}=\eta_{\mu\nu}+h_{\mu\nu},
  \qquad |h_{\mu\nu}|\ll1.
\]
Após uma escolha conveniente de gauge, as perturbações radiativas satisfazem,
no vácuo,
\[
  \Box h^{\rm TT}_{ij}=0.
\]
O sobrescrito TT indica gauge transversal e sem traço.

\section{Polarizações}

Uma onda propagando-se na direção $z$ possui duas polarizações independentes:
\[
  h^{\rm TT}_{ij}=
  \begin{pmatrix}
  h_+ & h_\times & 0\\
  h_\times & -h_+ & 0\\
  0&0&0
  \end{pmatrix}.
\]
Elas deformam distâncias transversais de modos característicos.

\section{Detecção}

Interferômetros como LIGO, Virgo e KAGRA medem diferenças relativas de
comprimento entre braços quase perpendiculares. O sinal adimensional típico é
a deformação
\[
  h\sim\frac{\Delta L}{L}.
\]
Na prática, detectar ondas gravitacionais exige modelagem de ruído, filtros
casados e comparação com famílias de formas de onda.

\section{Emissão}

A radiação gravitacional dominante vem de variações temporais do momento de
quadrupolo da fonte. Uma estimativa dimensional para a potência emitida é
\[
  P\sim \frac{G}{c^5}\left(\frac{\dd^3 Q}{\dd t^3}\right)^2,
\]
onde $Q$ representa uma escala de quadrupolo de massa. A presença de $c^5/G$
mostra por que a emissão é pequena para sistemas cotidianos e enorme para
binárias compactas.

\section{Gauge e graus de liberdade}

Uma perturbação simétrica $h_{\mu\nu}$ possui inicialmente dez componentes.
Nem todas representam física. Parte delas muda quando fazemos uma pequena
transformação de coordenadas. Essa liberdade é chamada liberdade de gauge. A
escolha TT remove componentes não radiativas e deixa apenas dois graus de
liberdade físicos, correspondentes às polarizações $+$ e $\times$.

Essa situação é análoga ao eletromagnetismo, onde o potencial $A_\mu$ possui
liberdade de gauge e nem todos os seus componentes são observáveis. O campo
físico está nas combinações invariantes apropriadas. Em ondas gravitacionais,
o observável direto é a variação relativa de distâncias entre massas teste,
relacionada à curvatura perturbativa.

\section{Efeito sobre partículas teste}

Considere partículas inicialmente em repouso relativo formando um pequeno anel
no plano transversal à propagação da onda. Uma polarização $+$ alonga o anel em
uma direção enquanto o comprime na direção perpendicular; meio período depois,
os papéis se trocam. A polarização $\times$ produz o mesmo padrão rotacionado
de $45^\circ$.

Esse efeito não deve ser imaginado como uma força empurrando partículas dentro
de um espaço fixo. A onda altera a geometria, e as partículas livres seguem
geodésicas nessa geometria oscilante. O detector mede uma diferença de tempo de
propagação da luz entre braços, que pode ser descrita como uma deformação
relativa $h$.

\section{Escalas de detecção}

Uma deformação típica detectável por interferômetros terrestres é
$h\sim10^{-21}$. Para braços de $4\,\mathrm{km}$, a variação efetiva de
comprimento é da ordem de
\[
  \Delta L\sim hL\sim 4\times10^{-18}\,\mathrm{m}.
\]
Esse número é muito menor que o tamanho de um próton. A detecção só é possível
porque o experimento mede fase acumulada da luz, usa cavidades ópticas para
aumentar o caminho efetivo e emprega análise estatística com formas de onda
teóricas.

\section{Binárias compactas}

Durante a fase inspiralante de uma binária compacta, a emissão de ondas
gravitacionais remove energia e momento angular orbital. A separação diminui,
a frequência orbital aumenta e o sinal observado apresenta um ``chirp''. A
frequência da onda gravitacional dominante é aproximadamente duas vezes a
frequência orbital, pois a emissão quadrupolar repete o padrão duas vezes por
órbita.

As últimas órbitas e a fusão exigem relatividade numérica para descrição
precisa. Depois da fusão, o buraco negro final relaxa por modos quase-normais,
produzindo o chamado ringdown. O sinal completo permite inferir massas, spins,
distância luminosidade e testes de consistência da relatividade geral.

\begin{activity}
Esboce qualitativamente amplitude e frequência de uma onda gravitacional de
binária compacta em função do tempo. Indique as fases inspiral, merger e
ringdown.
\end{activity}

\section{Problemas orientados}

\begin{enumerate}
  \item \textbf{Traço TT.} Verifique que a matriz de polarizações em gauge TT
  tem traço nulo e é transversal para propagação na direção $z$.

  \item \textbf{Anel de partículas.} Escreva as deformações relativas
  $\Delta x/L$ e $\Delta y/L$ para uma onda $+$ ideal. Desenhe quatro instantes
  separados por um quarto de período.

  \item \textbf{Estimativa de detector.} Para $h=10^{-21}$ e $L=4$ km, calcule
  $\Delta L$. Compare com escalas nucleares e discuta por que interferometria é
  necessária.

  \item \textbf{Ausência de monopolo.} Explique por que conservação de massa
  total impede radiação gravitacional monopolar em sistemas isolados.

  \item \textbf{Ausência de dipolo.} Explique por que conservação do centro de
  massa elimina radiação dipolar gravitacional dominante.

  \item \textbf{Chirp.} Use energia orbital newtoniana e emissão qualitativa
  de energia para argumentar por que a frequência orbital aumenta durante o
  inspiral.
\end{enumerate}

\section{Erros comuns e interpretação}

Um erro comum é pensar em ondas gravitacionais como ondas que se propagam
dentro do espaço, deformando uma espécie de meio material. A descrição correta
é mais sutil: a própria métrica oscila. Como a métrica determina tempos
próprios, distâncias e cones de luz, a passagem de uma onda altera relações de
medida entre massas teste. Essa alteração é pequena, transversal e possui duas
polarizações independentes.

Outro erro frequente é imaginar que qualquer movimento de massa produz ondas
gravitacionais intensas. Sistemas esfericamente simétricos não irradiam. Um
colapso perfeitamente esférico, por mais violento que seja, não emite ondas
gravitacionais. É a variação de quadrupolo que importa. Por isso binárias
compactas são fontes tão eficientes: elas têm grandes massas, velocidades
relativísticas e forte assimetria quadrupolar.

Também é importante distinguir detecção direta de evidência indireta. Antes de
2015, a perda orbital de pulsares binários já fornecia forte confirmação da
emissão de ondas gravitacionais. A detecção interferométrica acrescentou a
observação direta do sinal passando pela Terra, abrindo uma forma nova de
astronomia.

\clearpage
\section{Exercícios}

\begin{exercise}
Descreva qualitativamente como uma onda com polarização $+$ altera um anel de
partículas teste.
\end{exercise}

\begin{exercise}
Por que sistemas esfericamente simétricos não emitem ondas gravitacionais?
\end{exercise}

\chapter{Soluções esféricas, estrelas e Schwarzschild}

\begin{objetivos}
Aplicar as equações de Einstein a sistemas esfericamente simétricos e extrair
previsões observáveis da geometria de Schwarzschild.
\end{objetivos}

\section{Métrica de Schwarzschild}

A solução exterior esfericamente simétrica, estática e de vácuo é
\[
  \dd s^2=
  -\left(1-\frac{2GM}{r}\right)\dd t^2
  +\left(1-\frac{2GM}{r}\right)^{-1}\dd r^2
  +r^2\dd\Omega^2,
\]
com
\[
  \dd\Omega^2=\dd\theta^2+\sin^2\theta\,\dd\phi^2.
\]
Em unidades $G=c=1$, substituímos $GM$ por $M$.

\section{Redshift gravitacional}

Para observadores estáticos em Schwarzschild, a frequência de um fóton emitido
em $r_e$ e observado em $r_o$ satisfaz
\[
  \frac{\nu_o}{\nu_e}
  =
  \sqrt{
  \frac{1-2GM/r_e}{1-2GM/r_o}
  }.
\]
Se o observador está muito longe, $r_o\to\infty$, então
\[
  \nu_o=\nu_e\sqrt{1-\frac{2GM}{r_e}}.
\]

\section{Geodésicas e potencial efetivo}

No plano equatorial $\theta=\pi/2$, há duas constantes de movimento:
\[
  E=\left(1-\frac{2M}{r}\right)\dot t,
  \qquad
  L=r^2\dot\phi.
\]
Para partículas massivas,
\[
  \dot r^2+V_{\rm eff}(r)=E^2,
  \qquad
  V_{\rm eff}(r)=\left(1-\frac{2M}{r}\right)
  \left(1+\frac{L^2}{r^2}\right).
\]
Para fótons, o termo $1$ é substituído por $0$.

\section{Estrelas relativísticas}

Dentro de uma estrela estática e esfericamente simétrica, a matéria pode ser
modelada por fluido perfeito. A equação de equilíbrio relativístico é a
equação de Tolman-Oppenheimer-Volkoff:
\[
  \frac{\dd p}{\dd r}
  =
  -\frac{(\rho+p)(m+4\pi r^3p)}{r(r-2m)},
  \qquad
  \frac{\dd m}{\dd r}=4\pi r^2\rho.
\]
Ela reduz a equação hidrostática newtoniana quando $p\ll\rho$ e $2m/r\ll1$.

\section{Coordenada areal}

Na métrica de Schwarzschild, a coordenada $r$ não é distância radial própria.
Ela é definida pela área das esferas de simetria:
\[
  A=4\pi r^2.
\]
A distância radial própria entre $r_1$ e $r_2$ em uma fatia $t=\mathrm{const}$
é
\[
  \ell=\int_{r_1}^{r_2}
  \left(1-\frac{2M}{r}\right)^{-1/2}\dd r.
\]
Essa distinção é importante: perto do horizonte, a relação entre incremento de
coordenada radial e distância própria torna-se muito diferente da intuição
euclidiana.

\section{Órbitas circulares}

O potencial efetivo permite estudar órbitas circulares. Para partículas
massivas, uma órbita circular satisfaz
\[
  \dot r=0,\qquad \frac{\dd V_{\rm eff}}{\dd r}=0.
\]
A estabilidade exige ainda
\[
  \frac{\dd^2 V_{\rm eff}}{\dd r^2}>0.
\]
Em Schwarzschild, há órbitas circulares estáveis apenas para $r>6M$. O raio
$r=6M$ é a órbita circular estável mais interna, conhecida como ISCO. Em discos
de acreção ao redor de buracos negros, essa escala é fisicamente importante,
pois marca aproximadamente a borda interna de discos finos relativísticos.

\section{Órbita de fótons}

Para geodésicas nulas, o potencial efetivo é
\[
  V_{\rm eff}^{(\gamma)}=
  \left(1-\frac{2M}{r}\right)\frac{L^2}{r^2}.
\]
A órbita circular nula ocorre em $r=3M$. Essa órbita é instável: uma pequena
perturbação leva o fóton a escapar para longe ou cair em direção ao buraco
negro. A região próxima a $r=3M$ está ligada à formação da sombra de buracos
negros, pois raios luminosos podem executar trajetórias fortemente curvadas
antes de chegar ao observador.

\section{Testes clássicos}

Três previsões históricas tornaram a relatividade geral observacionalmente
marcante: o avanço do periélio de Mercúrio, o desvio da luz pelo Sol e o
redshift gravitacional. O avanço do periélio surge porque a equação radial
relativística contém correções que impedem a órbita de fechar exatamente. O
desvio da luz é maior que uma estimativa newtoniana ingênua porque tanto a
curvatura temporal quanto a espacial contribuem. O redshift gravitacional
segue da diferença entre ritmos de relógios estáticos em potenciais
diferentes.

\begin{example}
Para emissão em $r_e$ e observação no infinito, o fator de redshift em
Schwarzschild é
\[
  \frac{\nu_\infty}{\nu_e}=\sqrt{1-\frac{2M}{r_e}}.
\]
Se $r_e$ se aproxima de $2M$, a frequência observada no infinito tende a zero.
Isso não significa que o fóton ``perde toda energia localmente''; significa
que a comparação entre o emissor próximo ao horizonte e o observador no
infinito envolve uma diferença gravitacional extrema.
\end{example}

\section{Problemas orientados}

\begin{enumerate}
  \item \textbf{Raio de Schwarzschild.} Calcule $r_s=2GM/c^2$ para o Sol e para
  a Terra. Compare com seus raios físicos.

  \item \textbf{Redshift fraco.} Expanda
  $\sqrt{1-2M/r}$ para $M/r\ll1$ e mostre que recupera
  $\Delta\nu/\nu\approx\Delta\Phi$.

  \item \textbf{Órbita de fótons.} Derive a condição $r=3M$ a partir do
  potencial efetivo nulo.

  \item \textbf{ISCO.} Pesquise algebricamente as condições de circularidade e
  estabilidade do potencial massivo e identifique por que $r=6M$ é especial.

  \item \textbf{Queda radial.} Para $L=0$, escreva a equação radial de uma
  partícula que cai a partir do repouso no infinito. Interprete $E=1$.

  \item \textbf{Atraso de Shapiro.} Explique qualitativamente por que um raio
  de luz que passa perto do Sol demora mais que em geometria plana.
\end{enumerate}

\section{Erros comuns e interpretação}

O primeiro cuidado em Schwarzschild é não confundir a coordenada $r$ com
distância radial própria. O raio $r$ é definido pela área das esferas de
simetria. Essa escolha é extremamente conveniente, mas a distância própria
radial envolve o fator $(1-2M/r)^{-1/2}$. Muitos paradoxos aparentes perto do
horizonte nascem de tratar $r$ como se fosse uma régua euclidiana.

Outro ponto importante é que o tempo coordenado $t$ é adaptado ao observador
distante, não ao observador em queda. Quando uma partícula se aproxima de
$r=2M$, o tempo $t$ necessário para cruzar o horizonte diverge. O tempo próprio
da partícula, porém, permanece finito. A divergência de $t$ expressa uma
limitação das coordenadas de Schwarzschild para cobrir o horizonte, não uma
barreira física local para quem cai.

Por fim, órbitas em Schwarzschild devem ser lidas com potencial efetivo, mas
esse potencial é uma ferramenta, não uma energia potencial newtoniana literal.
Ele resume a equação radial usando constantes de movimento. Sua forma permite
visualizar pontos de retorno, órbitas circulares e estabilidade, mas a dinâmica
subjacente continua sendo geodésica em espaço-tempo curvo.

\clearpage
\section{Exercícios}

\begin{exercise}
Use a métrica de Schwarzschild para obter a aceleração radial newtoniana no
limite $r\gg 2M$ e velocidades pequenas.
\end{exercise}

\begin{exercise}
Determine o raio da órbita circular nula em Schwarzschild usando o potencial
efetivo de fótons.
\end{exercise}

\chapter{Buracos negros}

\begin{objetivos}
Separar singularidades físicas de singularidades de coordenadas, compreender o
horizonte de eventos e introduzir buracos negros astrofísicos.
\end{objetivos}

\section{Horizonte}

Em Schwarzschild, $r=2M$ é o raio do horizonte de eventos. A divergência de
$g_{rr}$ nessa coordenada não é uma singularidade física. Invariantes de
curvatura permanecem finitos em $r=2M$ e divergem em $r=0$.

Um invariante útil é o escalar de Kretschmann:
\[
  R_{\alpha\beta\gamma\delta}R^{\alpha\beta\gamma\delta}
  =\frac{48M^2}{r^6}.
\]
Ele é finito no horizonte e diverge na singularidade central.

\section{Coordenadas regulares}

Coordenadas de Eddington-Finkelstein mostram que raios de luz podem atravessar
o horizonte para dentro em tempo finito. A dificuldade em $r=2M$ nas
coordenadas de Schwarzschild é análoga a uma má escolha de mapa para descrever
uma região perfeitamente regular.

\section{Buracos negros em rotação}

Buracos negros astrofísicos são descritos, em boa aproximação, pela métrica de
Kerr, caracterizada por massa $M$ e momento angular $J$. A rotação produz
arrastamento de referenciais e uma ergorregião, onde observadores não podem
permanecer estáticos em relação ao infinito.

\section{Termodinâmica}

Classicamente, a área do horizonte não diminui. Com efeitos quânticos, buracos
negros possuem temperatura
\[
  T_H=\frac{\hbar c^3}{8\pi G M k_B}
\]
para o caso de Schwarzschild. A entropia é proporcional à área do horizonte.

\section{Horizon versus singularidade}

O horizonte de eventos é uma superfície causal: ele separa eventos que podem
enviar sinais ao infinito futuro daqueles que não podem. A singularidade é uma
região onde a descrição clássica prevê curvaturas divergentes e geodésicas
incompletas. Confundir os dois conceitos leva a interpretações erradas. Um
observador em queda livre atravessa o horizonte de um buraco negro grande sem
encontrar uma parede local. Já a aproximação da singularidade envolve marés
crescentes que destroem qualquer objeto físico.

Para um observador distante em coordenadas de Schwarzschild, a queda de um
objeto parece congelar perto de $r=2M$, com sua luz cada vez mais desviada para
o vermelho. Para o objeto em queda, a travessia do horizonte ocorre em tempo
próprio finito. As duas descrições usam noções diferentes de tempo e não se
contradizem.

\section{Diagramas causais}

Diagramas de Penrose compactificam infinitos para representar a estrutura
causal global. Linhas de luz são desenhadas a $45^\circ$, o que facilita ver
quais regiões podem influenciar outras. Para Schwarzschild, o diagrama mostra
claramente que, dentro do horizonte futuro, qualquer trajetória causal aponta
inevitavelmente para menores valores de $r$. Nesse sentido, a singularidade
fica no futuro de todos os observadores que cruzam o horizonte.

\section{Buracos negros reais}

Buracos negros astrofísicos não são exatamente Schwarzschild. Eles giram,
acrescem matéria, podem estar em binárias e são perturbados pelo ambiente.
Ainda assim, a solução de Schwarzschild é o laboratório conceitual mais simples
para entender horizonte, geodésicas e redshift. A solução de Kerr generaliza o
caso para rotação e introduz fenômenos novos, como ergosfera e arrastamento de
referenciais.

Observacionalmente, buracos negros são estudados por órbitas estelares,
emissão de discos de acreção, jatos relativísticos, sombras observadas por
interferometria de base muito longa e ondas gravitacionais de fusões. Cada
canal testa uma região diferente da geometria.

\section{Leis da mecânica de buracos negros}

Classicamente, a área do horizonte desempenha papel análogo à entropia: ela
não diminui em processos físicos gerais. A gravidade superficial é análoga à
temperatura, e a massa obedece uma relação semelhante à primeira lei da
termodinâmica. A descoberta da radiação Hawking tornou essa analogia literal
em teoria quântica de campos em espaço-tempo curvo.

\begin{activity}
Compare a dependência $T_H\propto1/M$ com a intuição térmica comum. Por que
buracos negros maiores são mais frios? O que isso sugere sobre a evaporação de
buracos negros pequenos?
\end{activity}

\section{Problemas orientados}

\begin{enumerate}
  \item \textbf{Kretschmann.} Use
  $K=48M^2/r^6$ para comparar a curvatura no horizonte de buracos negros de
  massas $M$ e $10M$.

  \item \textbf{Tempo próprio de queda.} Discuta por que um observador em queda
  livre cruza o horizonte em tempo próprio finito, embora a coordenada $t$ de
  Schwarzschild divirja.

  \item \textbf{Cone de luz.} Desenhe qualitativamente cones de luz fora,
  sobre e dentro do horizonte. Indique por que dentro do horizonte diminuir
  $r$ é inevitável.

  \item \textbf{Ergosfera.} Explique a diferença entre horizonte e ergosfera em
  Kerr. Por que a ergosfera permite extração de energia?

  \item \textbf{Temperatura.} Estime como a temperatura de Hawking muda quando
  a massa é multiplicada por $10^6$.

  \item \textbf{Entropia de área.} Discuta por que a entropia proporcional à
  área, e não ao volume, sugere uma física gravitacional profundamente
  diferente da termodinâmica comum.
\end{enumerate}

\section{Erros comuns e interpretação}

Buracos negros não são aspiradores cósmicos. Longe do objeto, o campo
gravitacional de um buraco negro de massa $M$ é semelhante ao de qualquer corpo
esférico com a mesma massa. Se o Sol fosse substituído por um buraco negro de
mesma massa, as órbitas planetárias externas não seriam subitamente sugadas. O
que muda é a existência de uma região compacta com horizonte.

Também é incorreto dizer que nada acontece no horizonte em todos os sentidos.
Para um observador em queda livre, um horizonte de buraco negro suficientemente
grande pode ser atravessado sem efeito local dramático. Para a estrutura causal
global, porém, o horizonte é decisivo: após cruzá-lo, nenhum sinal causal pode
alcançar o infinito futuro. A frase correta depende de qual pergunta está sendo
feita: experiência local ou causalidade global.

A singularidade clássica em $r=0$ não deve ser interpretada como um ponto
material comum. Ela indica incompletude da teoria clássica em regime extremo.
Espera-se que gravidade quântica modifique essa descrição, mas ainda não há
uma teoria completa testada experimentalmente para substituir a previsão
clássica.

\clearpage
\section{Exercícios}

\begin{exercise}
Calcule o raio de Schwarzschild do Sol e da Terra. Comente por que esses
valores não significam que tais objetos sejam buracos negros.
\end{exercise}

\begin{exercise}
Explique a diferença entre horizonte de eventos e singularidade.
\end{exercise}

\chapter{Cosmologia relativística}

\begin{objetivos}
Usar simetria cosmológica para obter a métrica FLRW, as equações de Friedmann e
as relações observacionais básicas da expansão do Universo.
\end{objetivos}

\section{Princípio cosmológico}

Em grandes escalas, o Universo é aproximadamente homogêneo e isotrópico. A
métrica compatível com essas simetrias é
\[
  \dd s^2=-\dd t^2+a(t)^2
  \left[
  \frac{\dd r^2}{1-kr^2}+r^2\dd\Omega^2
  \right],
\]
onde $a(t)$ é o fator de escala e $k=0,\pm1$ representa a curvatura espacial.

\section{Redshift cosmológico}

O comprimento de onda da luz cresce com o fator de escala:
\[
  1+z=\frac{a(t_0)}{a(t_e)}.
\]
Esse redshift não é simplesmente Doppler local; ele expressa a evolução da
geometria entre emissão e observação.

\section{Equações de Friedmann}

Para fluido perfeito homogêneo,
\[
  H^2=\left(\frac{\dot a}{a}\right)^2
  =
  \frac{8\pi G}{3}\rho-\frac{k}{a^2}+\frac{\Lambda}{3},
\]
\[
  \frac{\ddot a}{a}
  =
  -\frac{4\pi G}{3}(\rho+3p)+\frac{\Lambda}{3}.
\]
A conservação local fornece
\[
  \dot\rho+3H(\rho+p)=0.
\]
Com uma equação de estado $p=w\rho$, segue
\[
  \rho\propto a^{-3(1+w)}.
\]
Matéria fria tem $w=0$; radiação tem $w=1/3$; energia escura do tipo constante
cosmológica tem $w=-1$.

\section{Distâncias em cosmologia}

Em um Universo em expansão, a palavra distância precisa ser qualificada. A
distância própria entre duas galáxias em uma fatia de tempo cosmológico cresce
com $a(t)$. A distância comóvel remove essa expansão e permanece fixa para
galáxias que acompanham o fluxo de Hubble. A distância luminosidade é definida
pela relação entre luminosidade intrínseca e fluxo observado:
\[
  F=\frac{L}{4\pi d_L^2}.
\]
A distância angular relaciona tamanho físico e tamanho angular:
\[
  \theta=\frac{\ell}{d_A}.
\]
Em cosmologia relativística, $d_L$ e $d_A$ não são iguais; eles obedecem
\[
  d_L=(1+z)^2d_A.
\]
Essa relação é chamada dualidade de Etherington e depende da conservação do
número de fótons e da propagação em geodésicas nulas.

\section{Horizontes cosmológicos}

O horizonte de partículas indica a maior distância comóvel de onde a luz pôde
chegar desde o início da expansão. O horizonte de eventos cosmológico, quando
existe, indica regiões que nunca poderão nos enviar sinais, mesmo esperando
tempo infinito. Em um Universo dominado por constante cosmológica, a expansão
acelerada produz um horizonte de eventos cosmológico.

Esses horizontes são diferentes do horizonte de um buraco negro. O horizonte de
buraco negro é associado a uma região da qual não se escapa para o infinito. O
horizonte cosmológico depende da história global da expansão e do observador
cosmológico considerado.

\section{Eras de dominação}

Com $p=w\rho$, a densidade escala como
\[
  \rho\propto a^{-3(1+w)}.
\]
Assim, radiação decai como $a^{-4}$: três potências vêm da diluição do número
de fótons e uma potência adicional vem do redshift da energia de cada fóton.
Matéria fria decai como $a^{-3}$, apenas por diluição do volume. A densidade de
uma constante cosmológica permanece constante.

Essa diferença produz uma sequência natural. No Universo muito jovem, a
radiação domina. Depois, a matéria passa a dominar e permite crescimento mais
eficiente de estruturas. Em tempos recentes, a energia escura domina e acelera
a expansão.

\section{Parâmetros observacionais}

Define-se a densidade crítica
\[
  \rho_c=\frac{3H_0^2}{8\pi G}.
\]
Os parâmetros adimensionais
\[
  \Omega_i=\frac{\rho_i}{\rho_c}
\]
medem a contribuição atual de cada componente. Para $k=0$,
\[
  \Omega_m+\Omega_r+\Omega_\Lambda=1.
\]
Observações de supernovas, radiação cósmica de fundo, oscilações acústicas de
bárions e estrutura em larga escala indicam um Universo aproximadamente plano,
com matéria escura fria e energia escura dominante hoje.

\section{Problemas conceituais}

O modelo cosmológico padrão é extremamente bem-sucedido, mas deixa questões
abertas. A natureza da matéria escura não é conhecida. A energia escura pode
ser uma constante cosmológica, mas seu valor observado é extraordinariamente
pequeno em comparação com escalas naturais de teoria quântica de campos. A
inflação resolve problemas de horizonte e planura em muitos modelos, mas sua
realização microscópica permanece em aberto. Essas questões mostram que a
cosmologia é tanto aplicação da relatividade geral quanto fronteira de física
fundamental.

\section{Problemas orientados}

\begin{enumerate}
  \item \textbf{Redshift e escala.} Calcule $a(t_e)/a(t_0)$ para $z=1$, $z=3$ e
  $z=1100$. Interprete o último valor no contexto da radiação cósmica de fundo.

  \item \textbf{Escala da radiação.} Derive $\rho_r\propto a^{-4}$ separando
  diluição de número e redshift de energia.

  \item \textbf{Universo de matéria.} Para $k=0$ e $\Lambda=0$, resolva
  $H^2=8\pi G\rho/3$ com $\rho\propto a^{-3}$ e obtenha $a\propto t^{2/3}$.

  \item \textbf{Universo de radiação.} Repita o problema anterior para
  $\rho\propto a^{-4}$ e obtenha $a\propto t^{1/2}$.

  \item \textbf{Aceleração.} Use
  $\ddot a/a=-(4\pi G/3)(\rho+3p)+\Lambda/3$ para determinar quando a expansão
  acelera.

  \item \textbf{Distâncias.} Explique por que distância luminosidade e
  distância angular diferem em um Universo em expansão. Qual papel é
  desempenhado pelo fator $(1+z)^2$?
\end{enumerate}

\section{Erros comuns e interpretação}

O redshift cosmológico não deve ser reduzido simplesmente a um efeito Doppler
global. Para galáxias próximas, a lei de Hubble pode ser aproximada por uma
interpretação Doppler. Em grandes distâncias, porém, a luz atravessa um
espaço-tempo cuja escala muda durante a propagação. A relação
$1+z=a(t_0)/a(t_e)$ resume esse efeito geométrico.

Outro cuidado é com a frase ``o Universo expande para dentro de algo''. A
métrica FLRW descreve a evolução das distâncias internas entre observadores
comóveis. Ela não exige um espaço externo no qual o Universo esteja embutido.
A pergunta correta é como a distância própria entre pontos comóveis muda com o
tempo cosmológico.

A expansão acelerada também não significa que todas as distâncias aumentem em
todos os sistemas. Sistemas ligados, como átomos, planetas, estrelas e
galáxias virializadas, não acompanham a expansão de Hubble em sua escala
interna. A expansão cosmológica descreve a dinâmica média em grandes escalas,
onde a hipótese de homogeneidade e isotropia se torna adequada.

\section{Síntese conceitual do curso}

A relatividade geral começa com uma mudança de linguagem: de forças em um
espaço absoluto para geometria dinâmica do espaço-tempo. A relatividade
especial ensina que o intervalo substitui tempo e espaço absolutos. O cálculo
tensorial fornece a gramática para escrever leis independentes de coordenadas.
A curvatura codifica marés e torna precisa a ideia de gravidade geométrica.
As equações de Einstein dizem como matéria e energia determinam essa curvatura.

As soluções e aproximações estudadas no final do curso mostram a força da
teoria. Schwarzschild descreve o exterior de corpos esféricos e introduz
horizontes. Ondas gravitacionais mostram que a geometria possui graus de
liberdade radiativos. Cosmologia aplica as equações de campo ao Universo como
um todo. Em todos os casos, a pergunta central permanece a mesma: quais
quantidades são invariantes, quais dependem do observador e quais são apenas
artefatos de coordenadas?

\clearpage
\section{Exercícios}

\begin{exercise}
Mostre que, para matéria fria em Universo plano sem constante cosmológica,
$a(t)\propto t^{2/3}$.
\end{exercise}

\begin{exercise}
Mostre que, para radiação em Universo plano sem constante cosmológica,
$a(t)\propto t^{1/2}$.
\end{exercise}

\chapter{Projetos finais}

Os projetos abaixo foram pensados para duas ou três semanas de trabalho em
duplas. Cada projeto deve resultar em um relatório curto e uma apresentação de
15 minutos.

\section{Projeto 1: geodésicas em Schwarzschild}

Implementar numericamente as equações de geodésicas no plano equatorial da
métrica de Schwarzschild. Comparar órbitas ligadas, espalhamento e queda radial.
Entregar gráficos de $r(\lambda)$ e da trajetória no plano.

\section{Projeto 2: redshift gravitacional}

Derivar a expressão do redshift em uma métrica estacionária e aplicá-la a
Schwarzschild. Discutir o limite fraco e estimar o efeito na superfície da
Terra, do Sol e de uma estrela de nêutrons típica.

\section{Projeto 3: ondas gravitacionais}

Estudar uma forma de onda simplificada de binária compacta inspiralante.
Relacionar crescimento de frequência e amplitude à perda de energia orbital.
Comparar com a ideia de sirene padrão.

\section{Projeto 4: cosmologias simples}

Resolver numericamente as equações de Friedmann para componentes de matéria,
radiação e constante cosmológica. Plotar $a(t)$ e $H(z)$ para diferentes
parâmetros.

\appendix

\chapter{Truques com índices e tensores}
\label{ap:indices}

Este apêndice fornece um algoritmo de trabalho para
manipular expressões com índices sem medo. O objetivo não é substituir a
interpretação geométrica dos tensores, mas dar ferramentas rápidas para
calcular, simplificar e conferir sinais.

\section{A filosofia dos índices}

Índices são uma linguagem para controlar somas. Um índice repetido é um índice
mudo: ele é somado e seu nome não importa. Por exemplo,
\[
  A_iB_i=A_jB_j=A_kB_k.
\]
As três expressões representam o mesmo número. O nome do índice repetido é
apenas uma variável de soma. Já um índice livre não pode ser renomeado em
apenas um lado da equação. A expressão
\[
  A_i=B_i+C_i
\]
é uma igualdade entre componentes de vetores. A expressão
\[
  A_i=B_j+C_j
\]
não é uma igualdade vetorial bem escrita, pois o lado esquerdo tem índice livre
$i$ e o lado direito não.

\begin{resultado}
Regra de ouro: em cada termo de uma equação, os índices livres devem ser os
mesmos; índices repetidos são somados e podem ser renomeados.
\end{resultado}

Em relatividade, costuma-se somar apenas quando um índice aparece uma vez
acima e uma vez abaixo:
\[
  A_\mu B^\mu.
\]
Em cálculo vetorial euclidiano tridimensional, muitas vezes escrevemos todos
os índices embaixo, usando implicitamente a métrica euclidiana
$\delta_{ij}$ para levantar e abaixar índices. Ao traduzir identidades 3D para
relatividade, é preciso recolocar a métrica correta e respeitar a assinatura.

\section{Delta de Kronecker}

O delta de Kronecker é definido por
\[
  \delta_{ij}=
  \begin{cases}
  1, & i=j,\\
  0, & i\ne j.
  \end{cases}
\]
Ele funciona como identidade para índices:
\[
  \delta_{ij}V_j=V_i,
  \qquad
  \delta_{ij}A_{jk}=A_{ik}.
\]
Em três dimensões,
\[
  \delta_{ii}=3.
\]
Em quatro dimensões,
\[
  \delta^\mu{}_\mu=4.
\]
O valor é a dimensão do espaço de índices que está sendo somado.

\begin{example}
Se $A_{ij}$ é uma matriz, então
\[
  \delta_{ij}A_{ij}=A_{ii}
\]
é o traço de $A$. Em quatro dimensões relativísticas, a contração
$T^\mu{}_\mu$ é o traço tensorial; se os dois índices estiverem ambos embaixo,
é preciso usar a métrica para contrair:
\[
  T^\mu{}_\mu=g^{\mu\nu}T_{\mu\nu}.
\]
\end{example}

\section{Índices mudos e troca de nomes}

Muitas simplificações dependem de trocar nomes de índices mudos. Considere
\[
  S_{ij}A_{ij},
\]
com $S_{ij}=S_{ji}$ simétrico e $A_{ij}=-A_{ji}$ antissimétrico. Como $i$ e
$j$ são mudos, podemos trocar seus nomes:
\[
  S_{ij}A_{ij}=S_{ji}A_{ji}.
\]
Usando as simetrias,
\[
  S_{ji}A_{ji}=S_{ij}(-A_{ij})=-S_{ij}A_{ij}.
\]
Portanto
\[
  S_{ij}A_{ij}=0.
\]

Esse truque aparece o tempo todo. Toda contração completa entre uma parte
simétrica e uma parte antissimétrica se anula.

\begin{activity}
Mostre que $F_{\mu\nu}U^\mu U^\nu=0$ se $F_{\mu\nu}$ é antissimétrico. A pista
é notar que $U^\mu U^\nu$ é simétrico em $\mu,\nu$.
\end{activity}

\section{Símbolo de Levi-Civita}

Em três dimensões euclidianas, o símbolo de Levi-Civita $\epsilon_{ijk}$ é
totalmente antissimétrico e definido por
\[
  \epsilon_{123}=1.
\]
Trocar dois índices muda o sinal; repetir índices dá zero. Assim,
\[
  \epsilon_{213}=-1,
  \qquad
  \epsilon_{113}=0.
\]

O produto vetorial pode ser escrito como
\[
  (A\times B)_i=\epsilon_{ijk}A_jB_k.
\]
Essa fórmula é mais segura que o determinante mnemônico quando as expressões
ficam longas.

\section{A identidade fundamental}

A identidade mais útil em três dimensões é
\[
  \epsilon_{ijk}\epsilon_{klm}
  =
  \delta_{il}\delta_{jm}
  -
  \delta_{im}\delta_{jl}.
\]
Ela transforma produtos de dois símbolos $\epsilon$ em deltas. Depois disso,
os deltas eliminam índices.

\begin{example}
Vamos obter o produto vetorial triplo:
\[
  [A\times(B\times C)]_i
  =
  \epsilon_{ijk}A_j(B\times C)_k.
\]
Como
\[
  (B\times C)_k=\epsilon_{klm}B_lC_m,
\]
temos
\[
  [A\times(B\times C)]_i
  =
  \epsilon_{ijk}\epsilon_{klm}A_jB_lC_m.
\]
Usando a identidade fundamental,
\[
  [A\times(B\times C)]_i
  =
  (\delta_{il}\delta_{jm}-\delta_{im}\delta_{jl})A_jB_lC_m.
\]
Logo,
\[
  [A\times(B\times C)]_i
  =
  B_i(A_jC_j)-C_i(A_jB_j).
\]
Em notação vetorial,
\[
  A\times(B\times C)=(A\cdot C)B-(A\cdot B)C.
\]
\end{example}

\section{Derivadas como índices}

Em coordenadas cartesianas euclidianas, escrevemos
\[
  \partial_i=\frac{\partial}{\partial x^i}.
\]
Então
\[
  \nabla\cdot A=\partial_iA_i,
  \qquad
  (\nabla\times A)_i=\epsilon_{ijk}\partial_jA_k.
\]
As identidades usuais de cálculo vetorial viram consequências de simetria e
antissimetria.

\begin{example}
Para mostrar que $\nabla\times\nabla f=0$, escrevemos
\[
  (\nabla\times\nabla f)_i
  =
  \epsilon_{ijk}\partial_j\partial_k f.
\]
O fator $\epsilon_{ijk}$ é antissimétrico em $j,k$, enquanto
$\partial_j\partial_k f$ é simétrico em $j,k$ se as derivadas mistas comutam.
Logo a contração é zero.
\end{example}

\begin{example}
Para mostrar que $\nabla\cdot(\nabla\times A)=0$,
\[
  \nabla\cdot(\nabla\times A)
  =
  \partial_i\epsilon_{ijk}\partial_jA_k
  =
  \epsilon_{ijk}\partial_i\partial_jA_k.
\]
Novamente, $\epsilon_{ijk}$ é antissimétrico em $i,j$, enquanto
$\partial_i\partial_j$ é simétrico. Portanto,
\[
  \nabla\cdot(\nabla\times A)=0.
\]
\end{example}

\section{Rotacional do rotacional}

Uma identidade muito usada em eletromagnetismo é
\[
  \nabla\times(\nabla\times A)=\nabla(\nabla\cdot A)-\nabla^2A.
\]
Em componentes,
\[
  [\nabla\times(\nabla\times A)]_i
  =
  \epsilon_{ijk}\partial_j(\nabla\times A)_k
  =
  \epsilon_{ijk}\epsilon_{klm}\partial_j\partial_l A_m.
\]
Usando a identidade fundamental,
\[
  [\nabla\times(\nabla\times A)]_i
  =
  (\delta_{il}\delta_{jm}-\delta_{im}\delta_{jl})
  \partial_j\partial_l A_m.
\]
Assim,
\[
  [\nabla\times(\nabla\times A)]_i
  =
  \partial_i\partial_mA_m-\partial_j\partial_jA_i,
\]
isto é,
\[
  \nabla\times(\nabla\times A)=\nabla(\nabla\cdot A)-\nabla^2A.
\]

\section{Tradução para relatividade}

No espaço-tempo, as regras de índices continuam, mas alguns objetos mudam de
papel:

\begin{itemize}
  \item o delta $\delta^\mu{}_\nu$ continua sendo a identidade;
  \item a métrica $g_{\mu\nu}$ é quem levanta e abaixa índices;
  \item a contração correta usa um índice acima e outro abaixo;
  \item o símbolo de Levi-Civita em quatro dimensões tem quatro índices;
  \item em espaço-tempo curvo, derivadas parciais de tensores são substituídas
  por derivadas covariantes.
\end{itemize}

Por exemplo, a norma de um vetor não é $V_\mu V_\mu$, mas
\[
  V_\mu V^\mu=g_{\mu\nu}V^\mu V^\nu.
\]
Com assinatura $(-,+,+,+)$, um vetor temporal possui norma negativa. Esse sinal
é físico e não deve ser apagado por analogia com a geometria euclidiana.

\section{Checklist para não errar}

Antes de aceitar uma conta com índices, verifique:

\begin{enumerate}
  \item Todo índice livre aparece igual em todos os termos?
  \item Algum índice aparece três vezes no mesmo termo?
  \item Todo índice repetido está realmente sendo somado?
  \item Em relatividade, a contração tem um índice acima e um abaixo?
  \item Alguma simetria ou antissimetria torna o termo zero?
  \item Algum delta pode eliminar um índice?
  \item Dois símbolos $\epsilon$ podem ser trocados por deltas?
  \item A derivada usada deveria ser parcial ou covariante?
\end{enumerate}

Esse checklist é simples, mas resolve a maior parte dos erros algébricos em
relatividade geral introdutória.

\chapter{Resumo de fórmulas}

\section{Relatividade especial}
\[
  \dd s^2=-\dd t^2+\dd x^2+\dd y^2+\dd z^2,
  \qquad
  \gamma=(1-v^2)^{-1/2}.
\]
\[
  U^\mu=\frac{\dd x^\mu}{\dd\tau},
  \qquad
  p^\mu=mU^\mu,
  \qquad
  E^2=\bm p^2+m^2.
\]

\section{Geometria}
\[
  \dd s^2=g_{\mu\nu}\dd x^\mu\dd x^\nu.
\]
\[
  \Gamma^\rho_{\mu\nu}
  =\frac{1}{2}g^{\rho\sigma}
  (\partial_\mu g_{\nu\sigma}
  +\partial_\nu g_{\mu\sigma}
  -\partial_\sigma g_{\mu\nu}).
\]
\[
  R^\rho{}_{\sigma\mu\nu}
  =\partial_\mu\Gamma^\rho_{\nu\sigma}
  -\partial_\nu\Gamma^\rho_{\mu\sigma}
  +\Gamma^\rho_{\mu\lambda}\Gamma^\lambda_{\nu\sigma}
  -\Gamma^\rho_{\nu\lambda}\Gamma^\lambda_{\mu\sigma}.
\]
\[
  G_{\mu\nu}=R_{\mu\nu}-\frac{1}{2}Rg_{\mu\nu}.
\]

\section{Einstein}
\[
  G_{\mu\nu}+\Lambda g_{\mu\nu}=8\pi G T_{\mu\nu}.
\]

\section{Schwarzschild}
\[
  \dd s^2=
  -\left(1-\frac{2GM}{r}\right)\dd t^2
  +\left(1-\frac{2GM}{r}\right)^{-1}\dd r^2
  +r^2\dd\Omega^2.
\]

\section{Cosmologia}
\[
  H^2=\frac{8\pi G}{3}\rho-\frac{k}{a^2}+\frac{\Lambda}{3},
  \qquad
  \dot\rho+3H(\rho+p)=0.
\]

\chapter{Bibliografia comentada}

\begin{itemize}
  \item Bernard Schutz, \textit{A First Course in General Relativity}, 3ª ed.,
  Cambridge University Press, 2022. Referência principal do curso; excelente
  para primeira exposição, com forte ligação entre cálculo tensorial e
  aplicações astrofísicas.

  \item Sean Carroll, \textit{Spacetime and Geometry}. Tratamento mais
  geométrico e moderno, recomendado para aprofundamento após os módulos de
  curvatura e equações de Einstein.

  \item James Hartle, \textit{Gravity}. Muito bom para intuição física,
  problemas e aplicações.

  \item Robert Wald, \textit{General Relativity}. Referência avançada, indicada
  para estudantes que queiram seguir pesquisa em gravitação matemática,
  buracos negros ou teoria clássica de campos em espaço-tempo curvo.

  \item Misner, Thorne e Wheeler, \textit{Gravitation}. Enciclopédico,
  historicamente importante e cheio de intuição geométrica.
\end{itemize}

\end{document}
